% !TeX program = xelatex
\documentclass[UTF8,fontset=fandol,zihao=-4,a4paper]{ctexart}
\usepackage[a4paper,margin=24mm,headheight=15pt]{geometry}
\usepackage{amsmath,amssymb,amsthm,mathtools,esint,booktabs,longtable,array,enumitem,xcolor,xurl,fancyhdr,microtype}
\usepackage[unicode,colorlinks=true,linkcolor=blue!45!black,citecolor=blue!45!black,urlcolor=blue!45!black]{hyperref}
\hypersetup{pdftitle={丘成桐第26题：独立核验、归约定理与现状判定}}
\newtheoremstyle{cn}{8pt}{8pt}{\normalfont}{0pt}{\bfseries}{。}{.4em}{}
\theoremstyle{cn}
\newtheorem{theorem}{定理}[section]
\newtheorem{lemma}[theorem]{引理}
\newtheorem{proposition}[theorem]{命题}
\newtheorem{corollary}[theorem]{推论}
\newtheorem{definition}[theorem]{定义}
\newtheorem{remark}[theorem]{说明}
\newtheorem{problem}[theorem]{问题}
\renewcommand{\proofname}{证明}
\DeclareMathOperator{\Ric}{Ric}
\DeclareMathOperator{\Rm}{Rm}
\DeclareMathOperator{\Vol}{Vol}
\DeclareMathOperator{\diam}{diam}
\DeclareMathOperator{\tr}{tr}
\DeclareMathOperator{\Id}{Id}
\DeclareMathOperator{\Met}{Met}
\DeclareMathOperator{\Hess}{Hess}
\newcommand{\op}{\mathrm{op}}
\newcommand{\R}{\mathbb R}
\newcommand{\Z}{\mathbb Z}
\newcommand{\CP}{\mathbb{CP}}
\newcommand{\CPb}{\overline{\mathbb{CP}}{}^2}
\newcommand{\HH}{\mathbb H}
\newcommand{\status}[1]{\textsf{\small〔#1〕}}
\numberwithin{equation}{section}
\setlength{\parskip}{.35em}
\setlength{\emergencystretch}{3em}
\linespread{1.12}
\allowdisplaybreaks[2]
\setlist{nosep,topsep=4pt,leftmargin=2em}
\pagestyle{fancy}\fancyhf{}
\fancyhead[L]{\small Yau 第26题：核验、归约与现状}
\fancyhead[R]{\small 2026-09-17}
\fancyfoot[C]{\thepage}
\newcolumntype{L}[1]{>{\raggedright\arraybackslash}p{#1}}
\title{丘成桐第 26 题（Ricci 夹逼与 Einstein 度量）\\[2mm]\large 独立核验、归约定理与现状判定}
\author{基于《求是引擎 · qiushi-yau26》资料包的独立研究报告}
\date{2026 年 9 月 17 日}
\begin{document}
\maketitle
\begin{abstract}
本报告对上传的《qiushi-yau26》资料包（六篇研究稿、两份计算附录与复算程序）作了逐项独立核验，并在此基础上给出关于 Yau 第 26 题的一个完整的\emph{现状判定}：哪些结论已被严格证明，哪些问题仍然开放，以及原题的最终解决在逻辑上等价于（或蕴含）哪些公认的公开问题。

\textbf{判定。}工作命题 $Q_n$（存在仅依赖维数的 Ricci 夹逼区间，使满足夹逼的闭流形必承载某个 Einstein 度量）在 $n=2,3$ 成立；在 $n\ge4$ \emph{目前无法被证明或否定}。这不是资料不足造成的：本报告证明（定理~\ref{thm:R}），$Q_4$ 成立将推出 $\CPb\#\CPb$ 等一批具体四维流形承载 Einstein 度量（或负分支的对应结论），而否定 $Q_4$ 则（据我们所知）要求对承载正数量曲率度量的四维流形建立全新的 Einstein 不存在性障碍；两者都是四维几何中公认的公开问题。$n\ge5$ 时，否定 $Q_n$ 以“存在不承载 Einstein 度量的闭 $n$ 维流形”为前提，而这样的流形目前一个也不知道。\textbf{因此本报告不宣称原题已解决}；它给出的是经过核验的定理、新的推论与精确的归约。

\textbf{已严格成立并经独立核验的主要结论。}（i）固定 $A_1$ 轨形解消上的 log-Weyl 构造：$\delta_\sigma(g_a)\le C/|\log a|\to0$（定理~\ref{thm:A1}；本报告用独立的符号与高精度数值计算复核，球面悬挂情形实测 $L\,\delta\to1/3$，与理论首项一致）；（ii）固定 Dehn-twist 曲面丛上的 $\delta_-\to0$ 构造；（iii）精确 Einstein 度量的统一拓扑间隙与近临界序列的必要塌缩；（iv）Hitchin--Thorpe、Seiberg--Witten、BCG 熵三类障碍与紧夹逼之间的定量关系（HT 与固定超临界 SW 比例被紧夹逼排除；SW 边界带与熵等号型流形既未构造也未排除）；（v）正交传递 $T^2$ 连接器的定量障碍。

\textbf{本次工作的新增内容。}（a）推论~\ref{cor:football}：$e_+(S^2\times S^2)=e_+(\CPb\#\CPb)=0$，并且相应度量族 GH 收敛到 $S^4/\Z_2$；这对 Ozuch（据其 2021 年报告）提出的 $L^\infty$ 夹逼公开问题给出肯定回答，同时说明“近 Einstein $\Rightarrow$ 附近有 Einstein 度量”的稳定性形式在四维两个符号分支均不成立（推论~\ref{cor:stability}）。（b）命题~\ref{prop:morse}：用 Morse 论证把正交传递障碍推广为底曲面的完全分类。（c）定理~\ref{thm:R}（归约定理）与定理~\ref{thm:W}（条件定理：若存在具有两个渐近双曲端的 $T^3\times\R$ 上的 Einstein 度量，则负分支 $P_4^-$ 不成立）。（d）对资料包主要定理的审稿式核验结论与勘误建议（第~\ref{sec:verification}节）。
\end{abstract}

\paragraph{阅读约定。}每个结论后标注其状态：\status{已证}表示本报告给出完整证明或已逐步核验资料包中的完整证明；\status{引用}表示依赖所列文献的定理；\status{条件}表示依赖明确列出的未证假设；\status{开放}表示未解决。范数约定：$\delta_\sigma(g)=\sup_M\|g^{-1}\Ric_g-\sigma\Id\|_{\op}$，$e_\sigma(M)=\inf_g\delta_\sigma(g)$，$\sigma\in\{+1,-1\}$。

\tableofcontents
\section{问题、量词与分支归约}\label{sec:problem}

\subsection{题文与工作命题}
资料包 \cite{Pkg} 所据的题文（出自用户提供的重编本，第 26 题；原始出处为 \cite{Yau}）为：
\begin{quote}
Find for each $N$ constants $c_N<C_N$ so that if the Ricci curvature of a compact manifold satisfies $c_N\delta_{ij}\le R_{ij}\le C_N\delta_{ij}$ then the manifold admits an Einstein metric.
\end{quote}
（重编本的评注指向 Anderson \cite{Anderson89} 在附加几何假设下的近 Einstein 存在定理。）本次文献检索未能取得 1982 年原书（Yau, \emph{Problem section}, Ann. of Math. Studies 102）的页面，故与资料包一致，\emph{不宣称}上述文字与原页逐字相同；以下全部结论针对如下明确的数学表述。

全文中 $M$ 表示闭（紧、无边）、连通、光滑 $n$ 维流形，$g$ 为光滑 Riemann 度量；$ag\le\Ric_g\le bg$ 按二次型逐点理解。“$M$ 承载 Einstein 度量”指存在光滑度量 $h$ 与常数 $\lambda\in\R$ 使 $\Ric_h=\lambda h$；不要求 $h$ 接近 $g$，也不限制 $\lambda$ 的符号。

\begin{definition}
对 $n\ge2$ 与 $\sigma\in\{+1,-1\}$，定义
\begin{align*}
Q_n&:\ \exists a<b\ \ \forall (M^n,g):\ ag\le\Ric_g\le bg\ \Longrightarrow\ M\text{ 承载 Einstein 度量};\\
P_n^\sigma&:\ \exists\varepsilon\in(0,1)\ \ \forall (M^n,g):\ \delta_\sigma(g)\le\varepsilon\ \Longrightarrow\ M\text{ 承载 Einstein 度量}.
\end{align*}
\end{definition}

\begin{proposition}[分支归约]\status{已证}\label{prop:branches}
对每个 $n\ge2$，$Q_n\iff P_n^+\lor P_n^-$。
\end{proposition}
\begin{proof}
常数缩放 $\hat g=cg$（$c>0$）不改变 $(0,2)$ 型 Ricci 张量，故 $\hat g^{-1}\Ric_{\hat g}=c^{-1}g^{-1}\Ric_g$。若 $P_n^\sigma$ 以 $\varepsilon$ 成立，取 $[a,b]=[\sigma-\varepsilon,\sigma+\varepsilon]$ 即得 $Q_n$。反之设 $Q_n$ 对 $[a,b]$ 成立；任取 $[u,v]\subset[a,b]$，$u<v$，$uv>0$，令 $\mu=(u+v)/2$，$\sigma=\operatorname{sign}\mu$，$\varepsilon=(v-u)/|u+v|\in(0,1)$。若 $\delta_\sigma(h)\le\varepsilon$，则 $g=h/|\mu|$ 的 Ricci 特征值落在 $[u,v]$ 内，于是 $M$ 承载 Einstein 度量；这就是 $P_n^\sigma$。
\end{proof}

\begin{corollary}\label{cor:logic}
（1）证明 $Q_n$ 只需证明一个符号分支；（2）否定 $Q_n$ 必须在\emph{同一维数}同时否定 $P_n^+$ 与 $P_n^-$，两支可用不同流形；（3）否定 $P_n^\sigma$ 等价于存在闭流形序列 $(M_j^n,g_j)$，$\delta_\sigma(g_j)\to0$，且每个 $M_j$ 不承载\emph{任何}符号的 Einstein 度量。特别地，若某个不承载 Einstein 度量的固定流形满足 $e_\sigma(M)=0$，则 $P_n^\sigma$ 不成立。
\end{corollary}

\begin{remark}[含零区间与固定宽度]
若 $a<0<b$，任何度量放大后都满足夹逼，故含零内点的区间在存在无 Einstein 流形的维数中不可能有效。另一方面，Lohkamp 定理 \cite{Lohkamp} 给出仅依赖维数的 $a(n)>b(n)>0$，使\emph{每个} $n\ge3$ 维流形承载 $-a(n)g<\Ric<-b(n)g$ 的度量；因此四维中任何有效的负区间 $[-C,-c]$ 必须满足 $C/c<a(4)/b(4)$，而资料包还证明了 $a(4)/b(4)>3$（$S^1\times S^3$ 上严格 Hitchin--Thorpe 端点的推论）。这些只限制区间的宽度，不触及“任意紧夹逼”的原题。
\end{remark}

\section{已解决的维数：$n=2$ 与 $n=3$}\label{sec:lowdim}

\begin{theorem}\status{已证（依赖所引经典定理）}\label{thm:lowdim}
$Q_2$ 与 $Q_3$ 成立。更精确地：
\begin{enumerate}
\item $n=2$：每个闭曲面承载常曲率度量，故任意区间 $[a,b]$ 均可。
\item $n=3$，正分支：任意 $0<a<b$ 均可（$P_3^+$ 对一切 $\varepsilon<1$ 成立）。
\item $n=3$，负分支：$P_3^-$ 对一切 $\varepsilon<1/3$ 成立。
\end{enumerate}
\end{theorem}
\begin{proof}
（1）一致化定理；二维 Einstein 度量即常 Gauss 曲率度量。

（2）Hamilton 定理 \cite{Hamilton}：闭三维流形上 $\Ric>0$ 的度量沿（归一化）Ricci 流收敛到常正截面曲率度量，后者是 Einstein 度量。

（3）三维中 Weyl 张量为零，曲率算子在 Ricci 特征标架下对角，截面曲率为 $K_{ij}=\frac12(r_i+r_j-r_k)$。若 $r_i\in[-1-\varepsilon,-1+\varepsilon]$ 且 $\varepsilon<1/3$，则一切二平面的截面曲率 $\le\frac12(-1+3\varepsilon)<0$。闭的负截面曲率三维流形是非球面的、不可约的，其基本群不含 $\Z^2$（Preissmann），由几何化定理（Perelman \cite{Perelman}）它承载双曲度量，即 $\Ric=-2g_{\mathrm{hyp}}$ 的 Einstein 度量。
\end{proof}

三维的论证依赖“Ricci 控制全曲率”；从四维起 Weyl 张量不受 Ricci 夹逼的代数控制，这正是第~\ref{sec:logweyl}节构造所利用的自由度。

\section{四维的第一类构造：固定 $A_1$ 解消上的 log-Weyl 度量}\label{sec:logweyl}

本节给出资料包中 \texttt{thm:A1} 的自足陈述与证明，并记录本次独立核验。机制是：Ricci 线性化算子“看不见”Weyl 型二次项，因而可以在长度为 $L\sim|\log a|$ 的对数环带上以 $O(1/L)$ 的点态 Ricci 代价，把轨形节点处的 Weyl 曲率缓慢地换成 Eguchi--Hanson（EH）核心所要求的 Weyl 曲率。

\subsection{陈述}
设 $(M_0,g_0)$ 为闭四维 Einstein 轨形，$\Ric_{g_0}=\sigma g_0$，$\sigma\in\{\pm1\}$，奇点为有限个点 $p_1,\dots,p_m$，每点的局部群均为 $\{\pm1\}\subset O(4)$（$A_1$ 型）。EH 空间 $N=T^*S^2$ 的端部渐近于 $(\R^4\setminus B)/\{\pm1\}$；因 $-1$ 位于 $O(4)$ 的中心，任一 $A\in O(4)$ 都给出端部与轨形坐标卡之间的等距识别（\emph{标记}）。$O(4)$ 的两个连通分支对应两种\emph{手性}：相对于节点处给定的定向，EH 的零截面自交数分别为 $-2$ 与 $+2$。对每个节点选定标记后，把锥邻域换成 $N$ 的圆盘丛，得到光滑闭流形 $M$；其微分同胚型只依赖各节点的手性。

\begin{theorem}[log-Weyl 构造]\status{已证}\label{thm:A1}
对上述任意 $(M_0,g_0)$ 与任意手性选择，相应的光滑流形 $M$ 上存在光滑度量族 $g_a$（$0<a<a_0$），满足
\begin{equation}\label{eq:A1bound}
\delta_\sigma(g_a)\le\frac{C}{|\log a|}\xrightarrow[a\to0]{}0,
\end{equation}
$(M,g_a)$ Gromov--Hausdorff 收敛到 $(M_0,g_0)$，且各核心按 $a^{-2}$ 放大后光滑收敛到 EH 度量。常数 $C$ 依赖 $(M_0,g_0)$ 与标记，不依赖 $a$。特别地 $e_\sigma(M)=0$。
\end{theorem}

\subsection{三个引理}
取 EH 度量 $k=W^{-1}d\rho^2+\frac{\rho^2}{4}(\sigma_1^2+\sigma_2^2+W\sigma_3^2)$，$W=1-\rho^{-4}$，$\rho\ge1$，其中 $\sigma_i$ 为 $SU(2)$ 的左不变 1-形式，$\sigma_3=d\psi+\cos\theta\,d\phi$，$\psi$ 的周期为 $2\pi$。在端部的渐近坐标 $x$（$|x|=\rho$）中，欧氏度量 $e=d\rho^2+\frac{\rho^2}{4}(\sigma_1^2+\sigma_2^2+\sigma_3^2)$，$k=e+O(\rho^{-4})$，并有 $\langle Jx,dx\rangle=\frac{\rho^2}{2}\sigma_3$，$J$ 为使零截面全纯的复结构。

\begin{lemma}[显式核心修正]\status{已证}\label{lem:core}
令
\begin{align*}
\tilde h={}&-\frac{\rho^2(7\rho^2-5)}{36(\rho^2-1)(\rho^2+1)^2}\,d\rho^2-\frac{5(\rho^4-1)}{144}(\sigma_1^2+\sigma_2^2)
+\Bigl(-\frac{\rho^4}{72}+\frac1{16}-\frac{1}{12\rho^{2}}+\frac{5}{144\rho^{4}}\Bigr)\sigma_3^2,\\
H_*={}&-\frac5{36}\bigl(\rho^2e-(x\cdot dx)^2\bigr)+\frac1{12}\langle Jx,dx\rangle^2 .
\end{align*}
则 $\tilde h$ 是 $T^*S^2$ 上的光滑对称二张量，$D\Ric_k(\tilde h)=k$，且 $|\nabla_k^j(\tilde h-H_*)|_k\le C\rho^{-2-j}$、$|\nabla_k^j\tilde h|_k\le C(1+\rho)^{2-j}$（$j=0,1,2$）。
\end{lemma}
\begin{proof}
考虑双轴度量族 $q_\Lambda=F_\Lambda^{-1}d\rho^2+\frac{\rho^2}{4}(\sigma_1^2+\sigma_2^2+G_\Lambda\sigma_3^2)$，其中 $c_\Lambda=(1-\Lambda/4)^2-1$，$F_\Lambda=(1+c_\Lambda\rho^2)G_\Lambda$，$G_\Lambda$ 是
\[
\rho(1+c_\Lambda\rho^2)G'+(4+3c_\Lambda\rho^2)G=4-\Lambda\rho^2,\qquad G(1)=0
\]
的解。对一般双轴度量 $X\,d\rho^2+Y(\sigma_1^2+\sigma_2^2)+Z\sigma_3^2$ 直接计算 Ricci 张量（它在标架 $d\rho,\sigma_1,\sigma_2,\sigma_3$ 下对角），代入上式并用该 ODE 消去 $G',G''$，三个 Ricci 特征值恒等于 $\Lambda$，即 $\Ric(q_\Lambda)=\Lambda q_\Lambda$。在 bolt $\rho=1$ 处 $F'_\Lambda(1)=4-\Lambda$，$G'_\Lambda(1)=(4-\Lambda)/(1+c_\Lambda)$，故 $F'_\Lambda(1)G'_\Lambda(1)=16$ 对所有 $\Lambda$ 成立；这正是固定周期 $2\pi$ 下无锥角的条件，于是 $q_\Lambda$ 是同一光滑流形上关于 $\Lambda$ 光滑的光滑度量族，$q_0=k$。对 $\Ric(q_\Lambda)=\Lambda q_\Lambda$ 在 $\Lambda=0$ 求导得 $D\Ric_k(\partial_\Lambda q_\Lambda|_0)=k$。由 ODE 解出 $\partial_\Lambda G|_0=\frac{\rho^2}{12}+\frac1{4\rho^2}-\frac1{3\rho^4}$。取在 bolt 处为零的光滑向量场 $V=\frac5{72}(\rho^3-\rho^{-1})\partial_\rho$；因 $\Ric(k)=0$，$D\Ric_k(\mathcal L_Vk)=0$。直接计算表明 $\tilde h=\partial_\Lambda q_\Lambda|_0-\mathcal L_Vk$ 恰为所述公式，并在正交标架中
\[
\tilde h-H_*:\quad \frac{5-7\rho^2}{36\rho^2(\rho^2+1)},\quad \frac5{36\rho^2},\quad \frac{9\rho^4-12\rho^2+5}{36\rho^2(\rho^4-1)}
\]
（分别为 $\rho\rho$、$11=22$、$33$ 分量），从而得到衰减估计。全部代数由本报告的脚本 \texttt{v1\_core.py}、\texttt{v1b\_family.py}、\texttt{v1c\_htilde.py} 以计算机代数独立验证。
\end{proof}

\begin{lemma}[共同径向规范与 Weyl 型差]\status{已证}\label{lem:weylgauge}
记 $R^{(1)}_{ikjl}=\delta_{ij}\delta_{kl}-\delta_{il}\delta_{jk}$，$B_{ikjl}=J_{ij}J_{kl}-J_{il}J_{kj}+2J_{ik}J_{jl}$。则 $H_*=-\frac13R^*_{ikjl}x^kx^l\,dx^idx^j$，其中 $R^*=\frac5{12}R^{(1)}-\frac1{12}B$ 是 Ricci 缩并等于 $e$ 的代数曲率张量。因此，若 $g_0$ 在节点处等变正规坐标下写作 $g_0=e+H_0+E_4$，$H_0=-\frac13\Rm_{g_0}(p)_{ikjl}x^kx^l$，$|\partial^jE_4|\le Cr^{4-j}$，则对任意标记 $A\in O(4)$，
\[
K:=H_0-\sigma A^*H_*=-\tfrac13C_{ikjl}x^kx^l\,dx^idx^j,\qquad C:=\Rm_{g_0}(p)-\sigma A^*R^*,
\]
其中 $C$ 是 Ricci 缩并为零的代数曲率张量（代数 Weyl 张量）。从而
\[
\tr_eK=0,\qquad \operatorname{div}_eK=0,\qquad \Delta_eK=0,\qquad K(x,\cdot)=0 .
\]
\end{lemma}
\begin{proof}
$R^{(1)}_{ikjl}x^kx^l=r^2\delta_{ij}-x_ix_j$，$B_{ikjl}x^kx^l=3(Jx)_i(Jx)_j$（用 $x^{\mathsf T}Jx=0$），故 $-\frac13R^*xx=-\frac5{36}(r^2e-x\otimes x)+\frac1{12}Jx\otimes Jx=H_*$。$R^{(1)}$ 与 $B$ 的 Ricci 缩并均为 $3e$，所以 $\Ric(R^*)=(\frac{15}{12}-\frac3{12})e=e$。$g_0$ 的 Einstein 条件给 $\Ric(\Rm_{g_0}(p))=\sigma e$，故 $C$ 无迹。四个等式由 $C$ 的曲率对称性与无迹性逐项得到：$\partial_iK_{ij}=-\frac13(C_{iijl}x^l+C_{ikji}x^k)=0$（第一项由反对称性为零，第二项是 Ricci 型缩并），$\Delta K_{ij}=-\frac23C_{ikjk}=0$，$\tr K=-\frac13C_{ikil}x^kx^l=0$，$K_{ij}x^j=-\frac13C_{ikjl}x^kx^lx^j=0$。同样的逐项计算对任意代数曲率张量 $R$ 给出：$H_R:=-\frac13R_{ikjl}x^kx^l\,dx^idx^j$ 满足 $H_R(x,\cdot)=0$ 与
\begin{equation}\label{eq:DRicHR}
D\Ric_e(H_R)=\Ric(R)\quad(\text{常值张量}),
\end{equation}
这里 $\Delta H_R=-\frac23\Ric(R)$，$\partial_i\operatorname{div}(H_R)_j=\frac13\Ric(R)_{ij}$，$\nabla^2\tr H_R=-\frac23\Ric(R)$。特别地 $D\Ric_e(\sigma A^*H_*)=\sigma e$。$A_1$ 群作用 $x\mapsto-x$ 保持一切偶次张量，故上述对象均下降到商。
\end{proof}

\begin{remark}
$R^*$ 的自对偶 Weyl 部分为 $-\frac1{12}(B-R^{(1)})$，在 $\omega_J$ 上的特征值为 $-s/12$（$s=4$），恰使 $\mathbf R_+(\omega_J)=(W_++\frac s{12})\omega_J=0$；（约定：取使 $\omega_J$ 自对偶的定向，即 $J$ 的复定向，此时零截面自交数为 $-2$；曲率算子以单位球面为 $\Id$ 归一化；则 $W(R^*)=W^+(R^*)$，谱为 $(-\frac13,\frac16,\frac16)$。）这与 Biquard \cite{Biquard} 的一阶去奇异化障碍 $\mathbf R_+(\omega_J)=0$ 一致：核心修正只有在远场曲率满足该条件时才存在。log-Weyl 构造并不要求环境轨形满足该条件，差额由下一引理的慢插值吸收。
\end{remark}

\begin{lemma}[对数插值恒等式]\status{已证}\label{lem:logid}
设 $d\ge3$，$K$ 为 $\R^d$ 上满足引理~\ref{lem:weylgauge}中四个等式的二次齐次对称二张量，$\eta$ 为 $u=\log r$ 的任意光滑函数，则
\begin{equation}\label{eq:logid}
D\Ric_e\bigl[\eta(\log r)K\bigr]=-\frac{\eta''+(d+2)\eta'}{2r^2}\,K .
\end{equation}
\end{lemma}
\begin{proof}
$D\Ric_e(h)=\frac12(-\Delta h+\nabla\operatorname{div}h+(\nabla\operatorname{div}h)^{\mathsf T}-\nabla^2\tr h)$。$\eta K$ 无迹；$\operatorname{div}(\eta K)=\eta'r^{-2}K(x,\cdot)+\eta\operatorname{div}K=0$。故 $D\Ric_e(\eta K)=-\frac12\Delta(\eta K)=-\frac12\bigl[(\Delta\eta)K+2\partial_r\eta\,\partial_rK\bigr]$，其中 $\partial_rK=2K/r$，$\Delta\eta=(\eta''+(d-2)\eta')/r^2$，合并即得。（$d=4,5$ 时另由 \texttt{v2\_logidentity.py} 对随机 Weyl 张量符号验证。）
\end{proof}

\subsection{定理~\ref{thm:A1}的证明}
取 $a\ll b=a^{1/2}\ll R=a^{1/4}\ll1$，$L=\log\frac{R}{8b}=\frac14|\log a|-\log8$。在每个节点的等变正规坐标 $x$（$r=|x|$）中，记 $h=\sigma\tilde h$（于是 $D\Ric_kh=\sigma k$），$H_1=\sigma A^*H_*$，$K=H_0-H_1$。取定 $[0,1]$ 上的光滑台阶函数 $\psi$（在 $0$ 附近为 $0$，在 $1$ 附近为 $1$）。分区定义 $g_a$：

\textbf{（I）核心 $r\le b$。}$g_a=G_a:=a^2(k+a^2h)$，核心坐标 $y=x/a$。因 $a^2|h|_k\le Cb^2\to0$，$k+a^2h$ 一致正定。Taylor 展开并利用 $\Ric(k)=0$、$D\Ric_kh=\sigma k$：
\[
\Ric(k+a^2h)=\sigma a^2k+\mathcal Q,\qquad |\mathcal Q|_k\le Ca^4\bigl(|h||\nabla^2h|+|\nabla h|^2+|\Rm_k||h|^2\bigr)\le Ca^4(1+\rho)^2 .
\]
常数缩放不改变 $(0,2)$ 型 Ricci，升指标乘 $a^{-2}$，故 $\|G_a^{-1}\Ric(G_a)-\sigma\Id\|_{\op}\le Ca^2(1+\rho)^2\le C(a^2+r^2)\le Ca$。

\textbf{（II）内截断 $b\le r\le2b$。}在物理坐标中 $G_a=e+H_1+\mathcal T_a$，$|\partial^j\mathcal T_a|\le C(a^4r^{-4-j}+a^4r^{-2-j})$（EH 尾项与引理~\ref{lem:core}的衰减）。令 $g_a=e+H_1+(1-\chi)\mathcal T_a$，$\chi$ 为 $[b,2b]$ 上的标准截断。截断使 Ricci 改变 $O(a^4b^{-6}+a^4b^{-4})$，二次模型 $e+H_1$ 自身的非线性残差为 $O(b^2)$，合计 $O(a)$。

\textbf{（III）对数环带 $4b\le r\le R/2$}（$2b\le r\le4b$ 上取 $g_a=e+H_1$）。令 $\eta=\psi(\log(r/4b)/L)$，$g_a=e+q$，$q=H_1+\eta K$。由齐次性 $|q|\le Cr^2$，$|\partial q|\le Cr$，$|\partial^2q|\le C$，故 $\Ric(e+q)=D\Ric_e(q)+\mathcal N(q)$，$|\mathcal N(q)|\le C(|q||\partial^2q|+|\partial q|^2)\le Cr^2$。由式~\eqref{eq:DRicHR} 有 $D\Ric_e(H_1)=\sigma e$，再由式~\eqref{eq:logid}（$d=4$）：
\[
\bigl|\Ric(g_a)-\sigma g_a\bigr|\le\frac{|\eta''+6\eta'|}{2r^2}|K|+Cr^2\le C\bigl(L^{-1}+R^2\bigr).
\]

\textbf{（IV）外截断 $R/2\le r\le R$ 与外部。}此处 $\eta\equiv1$，$e+H_1+K=e+H_0$；令 $g_a=e+H_0+\chi'E_4$，截断代价 $O(R^2)$；$r\ge R$ 处 $g_a=g_0$。

各区的二次模型与 $e$ 一致可比，核心在 bolt 处光滑，故 $g_a$ 是 $M$ 上的光滑正定度量；节点有限，修改球在 $a$ 小时互不相交，Ricci 是局部算子。取最大值得 $\delta_\sigma(g_a)\le C(L^{-1}+a^{1/2})\le C'/|\log a|$。被修改区域的直径为 $O(R)$，把它压到节点给出失真 $O(R)$ 的对应，故 $(M,g_a)\to(M_0,g_0)$（GH）；在核心的固定紧集上 $a^{-2}G_a=k+a^2h\to k$。\qed

\subsection{独立数值核验：球面悬挂 $S^4/\Z_2$}\label{sec:football}
对圆球 $S^4$（$\Ric=g$）在 $\R^4\oplus\R$ 中按 $(x,t)\mapsto(-x,t)$ 作商得到的“橄榄球”轨形（两个 $A_1$ 节点，$\sigma=+1$），以及负曲率类比的局部模型（$\sigma=-1$，节点邻域为双曲锥），整个构造落在 $U(2)$ 不变的双轴类中：此时 $H_0=-\frac{\sigma}{9}(r^2e-(x\cdot dx)^2)$，$K=\frac{\sigma r^4}{144}(\sigma_1^2+\sigma_2^2-2\sigma_3^2)$。本报告用计算机代数得到 $g_a$ 在五个区域的精确表达式及其 Ricci 特征值的精确公式，再以 $3\,\mathrm{ex}+60$ 位十进制精度（$a=10^{-\mathrm{ex}}$）在每区 400 个对数等距点上取值（$\sigma=+1$ 另算了 $a=10^{-128}$；脚本 \texttt{v3b\_football.py}；台阶函数取 $\psi=f(t)/(f(t)+f(1-t))$，$f=e^{-1/t}$，$\max\psi'=2$）。结果见表~\ref{tab:football}：上确界总在对数环带内取得，$L\,\delta_\sigma(g_a)\to1/3=\max\psi'/6$，与资料包给出的纯 Weyl 插值主项 $|\eta'|/6$ 完全一致；其余各区的误差为 $O(a)$ 或 $O(R^2)$，与证明中的估计同阶。

\begin{table}[ht]\centering\small
\caption{$S^4/\Z_2$ 上 log-Weyl 度量的点态 Ricci 缺陷（采样上确界；$\sigma=\pm1$ 的数值在所列位数内相同）}\label{tab:football}
\begin{tabular}{cccccccc}
\toprule
$a$ & $L$ & $\delta_\sigma(g_a)$ & $L\,\delta_\sigma$ & 核心 & 内截断 & 对数环带 & 外截断\\\midrule
$10^{-8}$ & $2.526$ & $0.135117$ & $0.34127$ & $1.6\times10^{-9}$ & $1.8\times10^{-8}$ & $0.1351$ & $1.3\times10^{-5}$\\
$10^{-16}$ & $7.131$ & $0.046850$ & $0.33408$ & $1.6\times10^{-17}$ & $1.8\times10^{-16}$ & $0.04685$ & $1.3\times10^{-9}$\\
$10^{-32}$ & $16.341$ & $0.020407$ & $0.33347$ & $1.5\times10^{-33}$ & $1.8\times10^{-32}$ & $0.02041$ & $1.3\times10^{-17}$\\
$10^{-64}$ & $34.762$ & $0.0095899$ & $0.33336$ & $1.4\times10^{-65}$ & $1.8\times10^{-64}$ & $0.00959$ & $1.3\times10^{-33}$\\
$10^{-128}$ & $71.603$ & $0.0046554$ & $0.33334$ & $1.2\times10^{-129}$ & $1.8\times10^{-128}$ & $0.004655$ & $1.3\times10^{-65}$\\
\bottomrule
\end{tabular}
\end{table}

\subsection{推论：具体流形与 Ozuch 问题}
\begin{corollary}[球面悬挂的两种解消]\status{已证}\label{cor:football}
$S^4/\Z_2$ 的两个节点各有两种手性，所得光滑流形恰有两个微分同胚型：
\[
S^2\times S^2\qquad\text{与}\qquad \CPb\#\CPb\ \ (\cong\CP^2\#\CP^2\ \text{的反定向}).
\]
两者都满足 $e_+=0$，且实现 $e_+=0$ 的度量族 GH 收敛到 $(S^4/\Z_2,g_{\rm round})$。
\end{corollary}
\begin{proof}
$M=D_1\cup_\phi D_2$，$D_i$ 为 $S^2$ 上 Euler 数 $e_i\in\{\pm2\}$（相对 $M$ 的一个整体定向）的圆盘丛，沿 $\R P^3$ 粘合。$\R P^3$ 的保向微分同胚同痕于恒同（Bonahon \cite{Bonahon}；亦为广义 Smale 猜想的推论），故 $M$ 的定向微分同胚型只依赖 $(e_1,e_2)$。$M$ 单连通，$b_2=2$；Mayer--Vietoris 序列表明两个零截面张成 $H_2(M;\Z)$ 的指数 $2$ 子格，其上交叉形式为 $\operatorname{diag}(e_1,e_2)$。圆球与 EH 的 $U(2)$ 对称性含公共的 $T^2$，两种手性的粘合均可取为 $T^2$ 等变，故 $M$ 承载有效光滑 $T^2$ 作用；由 Orlik--Raymond 分类 \cite{OR}，$M$ 是 $S^4,\pm\CP^2,S^2\times S^2$ 的连通和，$b_2=2$ 时只能是 $S^2\times S^2$、$\CP^2\#\CPb$、$\pm(\CP^2\#\CP^2)$。若 $e_1e_2<0$，形式不定；$\operatorname{diag}(1,-1)$ 中无平方为 $\pm2$ 的向量（$x^2-y^2=\pm2$ 无整数解），故形式为偶的双曲形式，$M\cong S^2\times S^2$。若 $e_1e_2>0$，形式定号，$M\cong\pm(\CP^2\#\CP^2)$。关于轨形的保向等距（交换两极）与反向等距（过两极的反射）表明四种标记恰给出这两型。$e_+=0$ 与 GH 收敛由定理~\ref{thm:A1}。
\end{proof}

\begin{remark}[Kirby 演算的独立验证]
$M=D_{-2}\cup(\text{沿 }-2\text{ 框平凡结的子午线、框数 }f\text{ 的 2-柄})\cup(\text{4-柄})$；两种 Hopf 纤维化在 $\R P^3$ 中共用同一圆周而框数相差 $1$。边界为 $S^3$ 要求 $|-2f-1|=1$，即 $f\in\{0,-1\}$：$(-2,0)$ 给 $S^2\times S^2$；$(-2,-1)$ 吹落 $-1$ 分支后给 $\CPb\#\CPb$。终审中另以 $T^2$ 权重独立复核了推论~\ref{cor:models}(1) 的奇偶性结论（翻转相对两节点得 spin 流形 $\#3(S^2\times S^2)$，翻转相邻两节点得 $3\CP^2\#3\CPb$）。
\end{remark}

\begin{remark}[与 Ozuch 问题的关系]
Ozuch \cite{OzuchII} 证明了：$M=T^*S^2\#S^4/\Z_2\#T^*S^2$ 上存在度量 $g_t\to S^4/\Z_2$（GH），满足 $\Ric(g_t)\ge3g_t$ 与 $\|\Ric(g_t)-3g_t\|_{L^p}\to0$（任意 $p<\infty$），但不存在满足 $\Ric=3g$ 的 Einstein 度量序列收敛到该轨形；据其 2021 年报告幻灯片（其中把 $M$ 写作 $S^2\times S^2$），他明确提出公开问题：能否做到 $\|\Ric(g_t)-3g_t\|_{L^\infty}\to0$？推论~\ref{cor:football}（归一化换为 $\Ric\approx3g$ 只需缩放 $g_a/3$）对\emph{两种}手性的解消都给出\textbf{肯定回答}，速率为 $O(1/|\log a|)$；再作缩放 $(1-\delta_a)g_a$ 还可同时保持单侧条件 $\Ric\ge3g$。文献检索未发现已有工作给出该结论或使用“对数 Weyl 插值”（混合手性 EH 粘合此前出现在 Brendle--Kapouleas \cite{BK} 关于 $T^4/\Z_2$ 的工作中，目标与结论不同）；资料包的构造稿应属新结果，建议单独成文。
\end{remark}

\begin{corollary}[更多模型]\status{已证；表中拓扑识别引自资料包}\label{cor:models}
（1）$\mathcal O=(S^2\times S^2)/\langle(r_\pi,r_\pi)\rangle$（四个 $A_1$ 节点，乘积 Einstein 度量）的各手性解消满足 $e_+=0$；资料包以拟环面分类给出其微分同胚型：$\CP^2\#5\CPb$、$2\CP^2\#4\CPb$、$3\CP^2\#3\CPb$、$\#3(S^2\times S^2)$ 及其反定向。
（2）对超椭圆曲线 $C_g,C_h$（$g,h\ge2$），$\mathcal O_{g,h}=(C_g\times C_h)/\langle\iota_g\times\iota_h\rangle$ 有 $4(g+1)(h+1)$ 个 $A_1$ 节点和 $\Ric=-g$ 的乘积轨形度量；其任意手性解消满足 $e_-=0$。全部取复方向的标准解消已承载负 Kähler--Einstein 度量（经整体平滑 $w^2=fq+tF$ 与同时解消）；含反手性节点的混合解消是否承载 Einstein 度量未知。
（3）只具有孤立 $\R^4/\{\pm1\}$ 奇点的紧双曲四维轨形（存在性见引理~\ref{lem:hyporb}）的 EH 解消满足 $e_-=0$。
\end{corollary}

\begin{lemma}[所需双曲轨形的存在性]\status{已证}\label{lem:hyporb}
存在紧、可定向的双曲四维轨形，其奇点全为孤立点且局部群均为 $\{\pm1\}$。
\end{lemma}
\begin{proof}
取 $k=\mathbb Q(\sqrt2)$，$f=x_1^2+\cdots+x_4^2-\sqrt2x_0^2$。$f$ 在恒同嵌入下符号为 $(4,1)$，在共轭嵌入下正定，故 $\Gamma=O'(f,\mathcal O_k)$ 是 $\operatorname{Isom}(\HH^4)$ 中的余紧算术格。取剩余特征 $p\ge5$ 且使主同余子群 $\Gamma(I)$ 无挠、保向的素理想 $I$。点反射 $s=\operatorname{diag}(-1,-1,-1,-1,1)\in\Gamma$ 正规化 $\Gamma(I)$，令 $\Gamma'=\Gamma(I)\rtimes\langle s\rangle$。$\Gamma'$ 的挠元必形如 $\gamma s$ 且为对合；$\gamma s\equiv s\pmod I$ 给 $\tr(\gamma s)\equiv-3$。保持 $\HH^4$ 的对合的 $+1$ 特征空间含类时向量，维数 $m\ge1$，迹为 $2m-5$；$2(m-1)\in I$ 且 $p\ge5$ 迫使 $m=1$，即 $\gamma s$ 是点反射，不动点孤立，在切空间上作用为 $-\Id$（保向）。同一点至多是一个点反射的不动点，故各稳定化子恰为 $\Z_2$。$\HH^4/\Gamma'$ 即为所求。
\end{proof}

\begin{corollary}[稳定性形式不成立]\status{已证（依赖所引隔离定理）}\label{cor:stability}
设 $\sigma=\pm1$。不存在函数 $\eta(\varepsilon)\to0$ 使得：$\delta_\sigma(g)\le\varepsilon$ 蕴含存在闭 Einstein 四维流形 $(N,h)$，$\Ric_h=\sigma h$，$d_{GH}((M,g),(N,h))\le\eta(\varepsilon)$。事实上，对推论~\ref{cor:football}（$\sigma=+1$）与推论~\ref{cor:models}(3)（$\sigma=-1$）中的 $g_a$，存在 $\eta_0>0$，使 $a$ 充分小时一切这样的 $(N,h)$ 满足 $d_{GH}\ge\eta_0$。
\end{corollary}
\begin{proof}
Ozuch（\cite{OzuchCPAM}，Theorem 0.9）证明：至少含一个 $\R^4/\Z_2$ 奇点的球面或双曲四维轨形不是光滑 Einstein 度量（固定 Einstein 常数）的 GH 极限；Huang--Ozuch（\cite{HuangOzuch}，Theorem 1.1）进一步给出这类轨形在 $\Ric=\pm3g$ 的 Einstein 轨形中的 GH 隔离。而 $(M,g_a)$ GH 收敛到该轨形，三角不等式即得结论。
\end{proof}

注意推论~\ref{cor:stability}在 $S^2\times S^2$ 上成立，而该流形承载乘积 Einstein 度量：\emph{输入度量附近无精确解}与\emph{底层流形无精确解}是两个命题，原题问的是后者。

\section{四维的第二类构造：固定 Dehn-twist 曲面丛}\label{sec:bundle}

\begin{theorem}[资料包长稿第六章；本次逐式复核]\status{已证}\label{thm:ED}
设 $F=\Sigma_p$、$B=\Sigma_q$（$p,q\ge2$），$D$ 为 $F$ 上一条非分离简单闭曲线的 Dehn twist，$u:\pi_1(B)\twoheadrightarrow\Z$ 为本原满同态，$E_D\to B$ 为单值化 $\gamma\mapsto[D]^{u(\gamma)}$ 的光滑曲面丛。则在这个\emph{固定}的闭光滑四维流形上存在度量族 $G_\ell$，体积恒为 $V_0$，满足
\[
\delta_-(G_\ell)\le C/L\longrightarrow0,
\]
其中 $L\sim\log(1/\ell)$ 为基底双曲领圈的半宽。因此 $e_-(E_D)=0$。
\end{theorem}
\begin{proof}[证明概要与核验]
纤维领圈上取 $h_z=du^2+a(u)^2(ds+z\zeta'(u)du)^2$，$a=a_0\cosh u$；对每个 $z$ 它是双曲度量（$K=-a''/a=-1$，面积形式与 $z$ 无关），且 $h_1=D^*h_0$。在基底的长双曲领圈上令 $z$ 从 $0$ 缓慢变到 $1$，领圈外 $z$ 为常数，于是 $g_B+h_{z}$ 经 $\operatorname{id}\times D$ 粘合下降到固定的 $E_D$。Ricci 缺陷矩阵（含两个底--纤维混合分量）在正交标架下可显式写出，其算子范数恰为 $\sqrt{v^4/4+\alpha^2+\beta^2}$，三项分别为 $O(L^{-2})$、$O(L^{-1})$、$O(L^{-1})$。本次核验用两条独立路线（坐标 Christoffel 符号；正交标架 Koszul 公式）符号复现了完整的 $4\times4$ Ricci 矩阵，并对显式轮廓数值得到 $L\cdot\sup\|\cdot\|_{\op}\to0.985$，确认速率为 $c/L$。
\end{proof}

资料包进一步证明：同时退化纤维与基底领圈并取平衡轮廓，$\delta_-(G_\varepsilon)=\frac{2}{\pi^4}\varepsilon^2+O(\varepsilon^4)$；一次真实的二阶张量校正把上确界残差降到 $O(\varepsilon^3)$；在校正背景的整个共形类内 $\inf_u\delta_-(u^2\widetilde G_\varepsilon)=\Theta(\varepsilon^3)$，且这些共形类不含任何符号的 Einstein 度量。复核结论：首系数 $2/\pi^4$、两个抵消恒等式、三个标量方程的分解、Bianchi 相容性与数值 PDE 验证均成立；这些下界与排除\emph{仅限所述共形类}（原稿对此陈述准确），不涉及 $E_D$ 上其他度量。

\begin{remark}
$E_D$ 是否承载 Einstein 度量未知。它在两个定向下都是辛流形，$\tau=0$，$\chi=4(p-1)(q-1)$；基本群无限且非几乎交换，故任何 Einstein 度量必有 $\lambda<0$；在 $\Ric=-3g$ 规范下，LeBrun 的数量曲率估计 $\int s^2\ge32\pi^2c_1^2$ 给 $\Vol\ge\frac{4\pi^2}{9}\chi$（等号当且仅当度量为 Kähler--Einstein），而 Gauss--Bonnet 给 $\Vol\le\frac{4\pi^2}{3}\chi$。两者相容，不产生矛盾；更细的必要条件见资料包长稿“曲面丛候选的拓扑”一节。
\end{remark}

\section{已知的 Einstein 障碍与紧夹逼的相容性}\label{sec:obstructions}

分支反例需要同一个 $M$ 同时满足 $e_\sigma(M)=0$ 与“无任何 Einstein 度量”。据我们所知，四维中现有的不存在性证书只有四类：Hitchin--Thorpe（HT）、Seiberg--Witten（LeBrun 型曲率估计 \cite{LeBrun01,IL}）、单纯体积（Gromov、Kotschick）、体积熵（Besson--Courtois--Gallot \cite{BCG}、Sambusetti \cite{Sambusetti}、Di Cerbo \cite{DiCerbo}）。本节汇总它们与紧夹逼的定量关系；各条均已逐式复核（第~\ref{sec:verification}节）。

\subsection{Hitchin--Thorpe：被紧夹逼自动满足}
\begin{proposition}\status{已证}\label{prop:HT}
闭定向四维流形若承载 $\delta_\sigma(g)\le\varepsilon<\frac12$ 的度量（$\sigma=\pm1$），则 $2\chi(M)>3|\tau(M)|$，并且 $2\chi\pm3\tau\ge\frac{\Vol_g}{4\pi^2}\min\{\frac23(1-\varepsilon)^2,\frac23(1+\varepsilon)(1-2\varepsilon)\}$。
\end{proposition}
\begin{proof}
$2\chi\pm3\tau=\frac1{4\pi^2}\int(2|W^\pm|^2+\frac{s^2}{24}-\frac{|E|^2}2)$，而 $\frac{s^2}{24}-\frac{|E|^2}{2}=\frac{s^2}6-\frac{|\Ric|^2}2=:F$ 对每个 Ricci 特征值是凹函数，在特征值箱体 $[\sigma-\varepsilon,\sigma+\varepsilon]^4$ 上的最小值在顶点取得；按高端点个数 $j=0,\dots,4$ 计算五个顶点值，最小者为 $\frac23(1-\varepsilon)^2$（$j=0$）与 $\frac23(1+\varepsilon)(1-2\varepsilon)$（$j=1$）中较小的一个，当 $\varepsilon<\frac12$ 时为正。
\end{proof}
（资料包基础篇给出的阈值 $1/(1+\sqrt3)$ 有效但不锐；长稿已给出端点 $1/2$，两处应统一。）因此违反或达到 HT 等号的流形（如 $\#_5\CP^2$、$S^1\times S^3$）不可能提供紧夹逼输入。

\subsection{Seiberg--Witten 禁区的误差窗口}
\begin{proposition}[资料包 \texttt{FND:SWgap}]\status{已证（依赖 LeBrun \cite{LeBrun01}）}\label{prop:SW}
设 $X$ 为具非平凡 SW 不变量的极小一般型复曲面，$M=X\#k\CPb\#\ell(S^1\times S^3)$，$K=k+4\ell$，$c=c_1^2(X)$，$r=K/c-\frac13$。若 $c/3\le K<c$，则 $M$ 上\emph{每个}度量满足
\[
\delta_-(g)\ge\Phi(r)=\begin{cases}\sqrt{r/2},&0\le r\le2(2-\sqrt3)^2,\\ 2\sqrt r/(\sqrt6+\sqrt r),&2(2-\sqrt3)^2\le r<\frac23.\end{cases}
\]
特别地 $\delta_-(g)\le\delta$ 蕴含 $K/c\le\frac13+\max\{2\delta^2,6\delta^2/(2-\delta)^2\}$（$\delta\le2-\sqrt3$ 时即 $\frac13+2\delta^2$）。
\end{proposition}
于是固定的超临界比例 $K/c>1/3$ 给出正间隙。注意 $\Phi(0)=0$：LeBrun 定理在 $K=c/3$ 处是非严格的，故恰在临界比例上、以及（由下面的定理~\ref{thm:U}）在 $(\frac13-\gamma)c\le k<\frac c3$ 的“边界带”内，流形不承载 Einstein 度量，而本命题不给出 $e_-$ 的正下界。

\subsection{统一拓扑间隙与必要塌缩}
\begin{theorem}[资料包 \texttt{thm:exact}；新结果]\status{已证；审稿结论“正确，需小幅修订”}\label{thm:U}
存在与 $X,k$ 无关的 $\gamma>0$：对一切满足 $b_+>1$ 的极小一般型复曲面 $X$ 与 $k\ge(\frac13-\gamma)c_1^2(X)$，$X\#k\CPb$ 不承载任何 Einstein 度量。更一般地，若 $M$ 对每个度量有满足 $(a^+)^2\ge A$ 的 monopole 类且 $\chi(M)\le LA$，则任何 Einstein 度量迫使 $(2\chi+3\tau)/A\ge\frac23+\eta(L)$。
\end{theorem}
证明把 LeBrun 估计中的近等号转化为非负密度 $\mathcal D$ 的小平均（$\fint\mathcal D\le436\delta^2+675d$，$d=(\frac13-\frac kA)_+$），用 Fubini 选点、带缓冲的局部等距展开的重数传输引理，以及\emph{不需体积下界}的 Cheeger--Tian \cite{CT} 四维 Einstein $\varepsilon$-正则性取得光滑局部极限；极限满足 $\mathcal D\equiv0$，即 $\Ric=-g$ 且 $\operatorname{spec}W^+=(-\frac16,\frac1{12},\frac1{12})$，而 Derdziński \cite{Derd} 型论证表明这在任何开集上不可能。

\begin{theorem}[局部正密度；资料包 \texttt{thm:local}]\status{已证（依赖 Cheeger--Naber \cite{CN}）}\label{thm:local}
对任意 $K,r,v>0$ 存在 $\eta_0(K,r,v)>0$：完备定向四维流形若满足 $|\Ric|_{\op}\le K$、$\Vol B_r(p)\ge v$，则 $\fint_{B_{2r}(p)}\mathcal D\ge\eta_0$。
\end{theorem}

\begin{theorem}[必要塌缩]\status{已证}\label{thm:collapse}
设 $M_i=X_i\#k_i\CPb$ 如上，$\delta_-(g_i)\to0$ 且 $(\frac13-k_i/A_i)_+\to0$。则对每个固定 $r,v>0$，$\Vol\{p:\Vol B_r(p)\ge v\}/\Vol(M_i)\to0$，且 $\diam(M_i,g_i)\to\infty$。
\end{theorem}
由此，在临界比例上（$d=0$）或沿 $d_i\to0$ 的序列，近 Einstein 输入只能来自几乎处处塌缩的度量；对边界带内固定的 $d>0$，定理~\ref{thm:local}型的局部密度下界只给出“$v$-非塌缩点的体积比例 $\le C(r,v)(436\delta^2+675d)$”，并不排除 $e_-=0$。另一方面，定理~\ref{thm:A1}的构造不会落入这一范畴：对 Kähler--Einstein $A_1$ 轨形，反手性节点产生正自交球面，使 SW 不变量为零，而全标准手性的解消是极小曲面（$k=0$）。因此 SW 边界带是一类“全部 Einstein 不存在性已知、紧夹逼输入既未构造也未排除”的固定流形。若资料包所述的弱调和片假设 (SC)（有界 Ricci、小 $L^2$ 曲率能量的塌缩球上存在原度量的统一弱调和展开片）成立，则得到无条件的一致间隙 $\delta^2+d\ge c>0$，从而 $d<c$ 的那部分边界带被一致排除。(SC) 目前未证；Y.~Li \cite{Li} 的四维 Ricci 流平滑定理给出的是 $|\Rm|\le C/t$，不足以推出所需的原度量控制。

\subsection{熵障碍：度二 BCG 与图流形预算}
\begin{proposition}\status{已证（依赖 BCG）}\label{prop:BCG}
设 $N$ 为闭定向实双曲四维流形，则 $N\#N$ 不承载任何 Einstein 度量，且 $e_-(N\#N)\ge\varepsilon_\chi=\frac{1-\sqrt{1-1/\chi_N}}{1+\sqrt{1-1/\chi_N}}=\frac1{4\chi_N}+O(\chi_N^{-2})$。
\end{proposition}
间隙依赖 $\chi_N$；令 $\chi_N\to\infty$ 时下界趋零，但这既不构造近 Einstein 度量，也不是维数一致的界（资料包另对普通点处的短连接把手给出覆盖体积障碍）。

对 Di Cerbo \cite{DiCerbo} 的 extended graph 四维流形（含纯实双曲块，沿 $T^3$ 尖点截面以仿射映射粘合；该范畴内已知无任何 Einstein 度量，证明用到 Connell--Suárez-Serrato \cite{CSS} 的最小体积计算），$\Ric\approx-g$ 归一化下任何 $\delta_-(g)=\delta\le\frac13$ 的度量满足
\[
\frac{V_*}{(1+\delta)^2}\le\Vol_g\le\frac{V_*}{(1-\delta)^2},\qquad \fint|W|^2\le\frac83\delta,\qquad \|W\|_\infty\ge\frac{1-3\delta}{\sqrt3},
\]
$V_*=12\pi^2\chi(M)$（$|W|$ 为 $\Lambda^2$ 上的 Hilbert--Schmidt 范数）。这些是相容性条件而非矛盾：局部极限可以是真双曲几何。图流形是否满足 $e_-=0$ 归结为第~\ref{sec:reduction}节的问题~W。

\subsection{正交传递 $T^2$ 对称类的完全排除}
本小节采用 $\Ric\approx-3g$ 规范，$\mathcal E=\Ric+3g$，$\epsilon=\|g^{-1}\mathcal E\|_{\infty,\op}$（换到 $\delta_-$ 须除以 $3$）。设完备四维 $(X,g)$ 具有自由、等距、正交传递的 $T^2$ 作用，$g=\gamma_B+e^\sigma M_{ab}d\xi^ad\xi^b$，$\det M=1$，$U:B\to\HH^2$ 描述 $M$，$e=|dU|^2$。称 $B$ 的一个端为\emph{精确端}，若在紧集外 $\gamma_B=dt^2+e^{2t}dz^2$、$e^\sigma=c\,e^{2t}$、$U$ 为常值，其中 $t\to+\infty$（外扩端）或 $t\to-\infty$（尖点端）。

\begin{theorem}[资料包 \texttt{thm:OT}]\status{已证}\label{thm:OT}
若 $B\cong S^1\times\R$ 且两端均为精确外扩端，则 $12-8\epsilon\le\mathcal B(\epsilon):=\bigl[\frac{4\epsilon}{\sqrt{3+\epsilon}}+\sqrt{2(3+\epsilon)+\frac{20\epsilon^2}{3+\epsilon}}\bigr]^2$；故 $\epsilon<\epsilon_*\approx0.24625$ 时不存在这样的度量。
\end{theorem}
其证明的上界部分实际上给出：对任意具有限个精确端的完备底面，
\begin{equation}\label{eq:supe}
\sup_Be\le\mathcal B(\epsilon),
\end{equation}
因为论证只用到完备性、$\Ric^4_f\ge-(3+\epsilon)\gamma_B$ 给出的 $CD(-(3+\epsilon),4)$ 条件（及 Erbar--Kuwada--Sturm／Bakry--Gentil--Ledoux 的 Fisher 信息 Riccati 不等式）以及 $e$ 与张力配对项具紧支集。本报告用一个逐点引理取代原证明中的极小分隔面，从而得到对底面拓扑的完全分类。

\begin{lemma}[临界点引理]\status{已证；新}\label{lem:critical}
令 $\mu=e^{\sigma/2}$（轨道面积的平方根）。在 $\mu$ 的任一临界点 $p$，若 $\nabla^2\mu(p)$ 不是正定的，则 $e(p)\ge12-8\epsilon$。
\end{lemma}
\begin{proof}
纤维块的迹给 $\Delta\sigma+|\nabla\sigma|^2=6-\tr\mathcal E_F$，$|\tr\mathcal E_F|\le2\epsilon$；底面块 $\Ric_g|_B=K_B\gamma_B-\nabla^2\sigma-\frac12d\sigma^2-\frac12U^*h_{\HH^2}=-3\gamma_B+\mathcal E_B$ 的无迹部分给 $(\nabla^2\sigma)_0+\frac12(d\sigma^2)_0=-\frac12(U^*h)_0-(\mathcal E_B)_0$。由 $\nabla^2\mu=\frac\mu2\nabla^2\sigma+\frac\mu4d\sigma^2$ 得
\[
\Delta\mu=\mu\bigl(3-\tfrac12\tr\mathcal E_F-\tfrac14|\nabla\sigma|^2\bigr),\qquad(\nabla^2\mu)_0=-\tfrac\mu4(U^*h)_0-\tfrac\mu2(\mathcal E_B)_0 .
\]
在临界点 $\Delta\mu\ge(3-\epsilon)\mu$。设 $\nabla^2\mu(p)$ 的特征值 $\lambda_1\ge\lambda_2$，$\lambda_2\le0$，则 $\|(\nabla^2\mu)_0\|_{\op}=\frac{\lambda_1-\lambda_2}2\ge\frac{\lambda_1+\lambda_2}2\ge\frac{(3-\epsilon)\mu}2$。另一方面 $\|(U^*h)_0\|_{\op}\le e/2$，$\|(\mathcal E_B)_0\|_{\op}\le\epsilon$，故 $\frac{3-\epsilon}2\le\frac e8+\frac\epsilon2$，即 $e\ge12-8\epsilon$。
\end{proof}

\begin{proposition}[正交传递类的底面分类]\status{已证；新}\label{prop:morse}
设 $B$ 连通、拓扑有限，全部端为精确端，且 $\epsilon<\epsilon_*$。则 $B\cong\R^2$ 且唯一的端是外扩端（AdS soliton 型），或 $B\cong S^1\times\R$ 且恰有一个外扩端与一个尖点端（形变尖点型）。特别地：不存在两端外扩的连接器（即定理~\ref{thm:OT}），也不存在底面亏格 $\ge1$ 或端数 $\ge3$ 的这类度量，从而在“端部精确、作用自由”的假设下，Di Cerbo 范畴中的 $\Sigma\times T^2$ 型块与 $T^3$ 连接颈均不能用正交传递度量实现 $\epsilon<\epsilon_*$。
\end{proposition}
\begin{proof}
由式~\eqref{eq:supe}与 $\mathcal B(\epsilon)<12-8\epsilon$，引理~\ref{lem:critical}表明 $\mu$ 的每个临界点处 Hessian 正定，即全为非退化极小点。把各端在精确区域内截断得紧曲面 $B_T$；外扩端边界上 $\nabla\mu$ 指向外，尖点端边界上指向内。带边 Morse 理论给 $\chi(B_T)-\chi(\partial_-B_T)=\#\{\text{极小点}\}$，而 $\partial_-B_T$ 为圆周之并，故 $\chi(B)=\#\{\text{极小点}\}\ge0$。$B$ 非紧（紧时 $\mu$ 有极大点，与临界点处 $\Delta\mu>0$ 矛盾），于是（$B$ 可定向时）$(\text{亏格},\text{端数})\in\{(0,1),(0,2)\}$；$B$ 不可定向时 $\chi\ge0$ 只允许开 Möbius 带，此时 $\mu$ 无临界点，而单个外扩端（尖点端）迫使内部最小（最大）点，矛盾。情形 $(0,1)$：若唯一的端是尖点端，则 $\mu$ 在内部取得最大值，矛盾。情形 $(0,2)$：$\mu$ 无临界点；两端皆外扩则 $\mu$ 有内部最小点，皆为尖点则有内部最大点，均矛盾。
\end{proof}

带联络扭转（非正交传递）的 $T^2$ 对称类不在上述结论内：此时式~\eqref{eq:supe}的 Bochner 论证中出现无定号项（资料包已指出），而引理~\ref{lem:critical}加强为 $e\ge12+2Q-8\epsilon$（$Q$ 为扭转能量密度）。资料包另证明了固定螺旋（Bianchi VII$_0$）标架类在反射中心后的不变区域障碍。

\section{归约定理：原题与公开问题的逻辑关系}\label{sec:reduction}

\subsection{$Q_4$ 成立意味着什么，否定 $Q_4$ 需要什么}
\begin{theorem}[归约定理]\status{已证}\label{thm:R}
\begin{enumerate}
\item[(a)] 若 $P_4^+$ 成立，则 $\CPb\#\CPb$（等价地 $\CP^2\#\CP^2$）承载 Einstein 度量；同样，$2\CP^2\#4\CPb$、$3\CP^2\#3\CPb$、$\#3(S^2\times S^2)$ 均承载 Einstein 度量（后三者作为 $(S^2\times S^2)/\Z_2$ 的解消，其拓扑识别引自资料包）。
\item[(b)] 若 $P_4^-$ 成立，则每个闭的负 Einstein 四维 $A_1$ 轨形的\emph{每一种}手性解消都承载 Einstein 度量（例如 $\mathcal O_{g,h}=(C_g\times C_h)/\Z_2$ 的全部混合手性解消），并且每个 Dehn-twist 曲面丛 $E_D$ 都承载 Einstein 度量。
\item[(c)] 因此 $Q_4$ 成立蕴含 (a) 的结论或 (b) 的结论。
\item[(d)] 若 $P_4^+$ 不成立，则存在闭四维流形 $M$：基本群有限，承载 $\delta_+<\frac12$ 的度量（从而承载正数量曲率度量，且 $M$ 及其一切有限覆盖满足严格不等式 $2\chi>3|\tau|$），而不承载任何 Einstein 度量。对这样的 $M$（不可定向时对其定向二重覆盖）及其一切有限覆盖：单纯体积与最小体积熵为零；任一 monopole 类满足 $a^+=0$，LeBrun 型估计的拓扑下界 $(a^+)^2>0$ 无从成立；HT 障碍失效。
\item[(e)] 对 $n\ge5$，否定 $Q_n$ 蕴含存在不承载任何 Einstein 度量的闭 $n$ 维流形。
\end{enumerate}
\end{theorem}
\begin{proof}
(a)(b)：推论~\ref{cor:football}、\ref{cor:models}与定理~\ref{thm:ED}给出所列流形满足 $e_\sigma=0$；$P_4^\sigma$ 的定义即给结论。(c) 由命题~\ref{prop:branches}。(d)：否定 $P_4^+$ 给出无 Einstein 度量的 $M$ 与 $\delta_+(g)<\frac12$ 的度量；Myers 定理给有限基本群，$s>0$；命题~\ref{prop:HT}给严格 HT，而提升度量满足同样的夹逼，故对有限覆盖亦然。有限基本群使单纯体积与体积熵为零，到双曲流形的映射度数为零，故 Gromov--Kotschick 与 BCG／Sambusetti 型证书不适用。正数量曲率度量下 Seiberg--Witten 方程只有可约解，故任一 monopole 类满足 $a^+_g=0$、$a^2\le0$；$b_+=1$ 时虚维数非负要求 $a^2\ge2\chi+3\tau>0$（严格 HT），零扰动下无墙，正数量曲率度量给出不变量为零，反定向同理。于是 LeBrun／Ishida--LeBrun \cite{LeBrun01,IL} 估计所需的 $(a^+)^2\ge a^2>0$ 无从成立。(e) 是定义的直接推论。
\end{proof}

\begin{remark}[这些结论为何构成“困难性证明”]
（1）据我们所知，(a) 中各流形是否承载 Einstein 度量是四维几何中长期公开的问题：已知的单连通正 Einstein 四维流形的例子只有（见 \cite{Viaclovsky,LeBrun03}） $S^4$、$\CP^2$、$S^2\times S^2$ 与 $\CP^2\#k\CPb$（$1\le k\le8$）；LeBrun 定理 \cite{LeBrun09} 表明承载辛结构的四维流形具有 $\lambda>0$ 的 Einstein 度量当且仅当它是 del Pezzo 曲面，而 $\CP^2\#\CP^2$ 不是近复流形，不在该定理范围内。对 (b) 中的混合手性解消，正自交球面使 SW 不变量在两个定向下都为零，同样没有已知的存在性或不存在性结果。（2）(d) 表明，据我们所知，文献中现有的全部 Einstein 不存在性证书对否定 $P_4^+$ 所需的流形均不适用；“正数量曲率、严格 HT、无 Einstein 度量”的四维流形目前一个也不知道；(e) 所要求的高维无 Einstein 流形也一个都不知道（$n\ge5$ 时尚无任何已知的 Einstein 存在性障碍）。（3）因此：\emph{证明} $Q_4$ 至少与在 $\CP^2\#\CP^2$（或相应负分支流形）上构造 Einstein 度量一样难；\emph{否定} $Q_4$ 至少与发现一种适用于正数量曲率四维流形的新 Einstein 障碍一样难，并且还需同时否定负分支。
\end{remark}

\subsection{负分支的一条可归约路线：问题 W}
“全部 Einstein 不存在性已被证明、而紧夹逼输入既未构造也未排除”的具体范畴目前有两类：SW 边界带 $(\frac13-\gamma)c\le k\le\frac c3$ 上的吹起流形（第~\ref{sec:obstructions}节；临界比例处必要塌缩），以及熵等号型流形，其代表是 Di Cerbo 的图流形。对后者，下述条件定理把问题归结为一个适定的渐近双曲 Einstein 存在性问题。

称完备 Einstein 流形 $(X^4,h)$（$\Ric_h=-3h$）的一个端为\emph{渐近双曲（AH）端}，若存在闭三维流形 $Y$ 与端部坐标使 $h=dt^2+e^{2t}\bar q(t)$，$t\ge t_0$，其中 $\bar q(t)$ 是 $Y$ 上的度量族，$\bar q(t)\to\gamma$，且 $\|\partial_t^j(\bar q(t)-\gamma)\|_{C^{2-j}(Y)}\le Ce^{-2t}$（$j=0,1,2$）。具光滑共形无穷远的共形紧 Einstein 度量满足此条件（Fefferman--Graham 展开）。

\begin{problem}[问题 W]\status{开放}\label{prob:W}
$T^3\times\R$ 上是否存在两端均为 AH 端的完备 Einstein 度量？
\end{problem}

\begin{theorem}[条件定理]\status{已证}\label{thm:W}
若问题~W 的答案是肯定的，则 $P_4^-$ 不成立：存在固定的闭四维流形 $M$（Di Cerbo 意义下由两个纯实双曲块组成的 extended graph 流形），$e_-(M)=0$ 而 $M$ 不承载任何 Einstein 度量。
\end{theorem}
\begin{proof}
设 $(X,h)$ 为所述度量。取两个具有环面尖点的有限体积实双曲四维流形 $N_\pm$（为简单计设各只有一个尖点；多尖点时逐对处理），其尖点为 $ds^2+e^{-2s}g^\pm_{\rm cusp}$。经仿射识别 $T^3\to$ 尖点截面并平移 $s$，可把尖点的一段写成 $dt^2+e^{2t}\gamma^\flat_\pm$，$\gamma^\flat_\pm$ 为 $T^3$ 上的平坦度量。在 $X$ 的每个端上取 $r_0\ge t_0$、$K\ge1$ 与光滑台阶函数 $\beta$，令
\[
\tilde q(t)=\bigl(1-\beta(\tfrac{t-r_0}{K})\bigr)\bar q(t)+\beta(\tfrac{t-r_0}{K})\gamma^\flat,\qquad g=dt^2+e^{2t}\tilde q(t).
\]
记 $J=\tilde q^{-1}\partial_t\tilde q$，则 $\|J\|_{C^1}+\|\partial_tJ\|_{C^0}\le C(K^{-1}+e^{-2r_0})$。Gaussian 坐标下的 Ricci 公式（资料包 \texttt{eq:fullRic}，已复核）给出 $\mathcal E=\Ric_g+3g$ 的三个块：
\begin{align*}
\mathcal E_{tt}&=-\tr J-\tfrac12\partial_t\tr J-\tfrac14\tr J^2,\qquad
\mathcal E_{t\cdot}=\tfrac12(\operatorname{div}_{\tilde q}J-d\tr J),\\
\mathcal E_T^\sharp&=e^{-2t}\Ric(\tilde q)^\sharp-\tfrac12\partial_tJ-\tfrac32J-\tfrac12(\tr J)\Id-\tfrac14(\tr J)J,
\end{align*}
混合块的物理范数还须乘 $e^{-t}$。$t\le r_0$ 时 $g=h$ 精确 Einstein；$t\ge r_0+K$ 时 $g$ 为精确双曲尖点；中间 $\|g^{-1}\mathcal E\|_{\op}\le C(K^{-1}+e^{-2r_0})$，$C$ 只依赖 $\bar q$ 与 $\gamma^\flat$ 的 $C^2$ 界。令 $K,r_0\to\infty$ 得到 $X$ 上缺陷任意小、两端在紧集外分别等于 $dt^2+e^{2t}\gamma^\flat_\pm$ 的完备度量；把两端与 $N_\pm$ 的尖点沿等距的一段粘合，得到闭流形 $M=\overline N_+\cup_{T^3}\overline N_-$ 上的光滑度量，外部为精确双曲。$M$ 的微分同胚型与参数无关，粘合映射是仿射的，故 $M$ 属于 Di Cerbo 定理的范畴，不承载任何 Einstein 度量；归一化到 $\Ric\approx-g$ 后 $e_-(M)=0$。
\end{proof}

\begin{remark}[关于问题 W 的已知约束]\label{rem:W}
（i）Cai--Galloway 定理 \cite{CG}：$\Ric\ge-3g$ 且衰减充分快的共形紧流形，只要共形无穷远有一个分支含非负数量曲率的代表，无穷远就连通。$T^3$ 上的共形类要么含平坦度量，要么是 Yamabe 负的；故问题 W 中两端的共形无穷远都必须是\emph{不共形于平坦度量}（Yamabe 负）的环面，平坦环面边界被排除。（ii）正交传递 $T^2$ 对称的解不存在：精确 Einstein 时 Bochner 公式给 $\frac12Le\ge\frac12e^2-3e$，AH 端上 $e\to0$，故 $\sup e\le6$；而环面面积函数 $\mu\to\infty$ 于两端，由引理~\ref{lem:critical}（$\epsilon=0$）与 Morse 计数必有 $e\ge12$ 的临界点，矛盾。（iii）固定螺旋标架的余齐性一类同样被排除（资料包 \texttt{prop:helical}）。（iv）反过来，Riccati 比较表明缺陷为 $\epsilon$ 的连接器长度 $\gtrsim\log(1/\epsilon)$；启发式地，非塌缩、曲率有界的近 Einstein 连接器序列（$\epsilon\to0$）以喉部为基点的极限应是 $T^3\times\R$ 型的完备 Einstein 度量，即问题 W 的解（端部的 AH 性质尚需论证）。在这个意义上，图流形路线与问题 W 密切相关。（v）资料包的全矩阵最小二乘试验与本次的 Gaussian 规范演化试验（第~\ref{sec:verification}节）都因 Cauchy 问题的椭圆不适定性而不能给出结论；正确的数值表述是以两端共形无穷远为 Dirichlet 数据的椭圆边值问题（Einstein--DeTurck 规范）。
\end{remark}

\section{高维 $n\ge5$}\label{sec:highdim}

\begin{proposition}\status{已证}\label{prop:product}
设 $e_\sigma(M^4)=0$，$(Y^k,g_Y)$ 为闭 Einstein 流形且 $\Ric_{g_Y}=\sigma g_Y$。则 $e_\sigma(M\times Y)=0$。从而对每个 $n\ge6$ 与每个符号，$P_n^\sigma$ 蕴含一批具体的闭 $n$ 维流形承载 Einstein 度量（本报告未查到关于它们的现成存在性结论），例如 $(\CPb\#\CPb)\times S^{n-4}$（$\sigma=+1$，球面取 $\Ric=g$ 的尺度）与 $M_-\times Y_{\rm hyp}^{\,n-4}$（$\sigma=-1$，$M_-$ 为推论~\ref{cor:models}(2) 的混合解消或 $E_D$，$Y_{\rm hyp}$ 为适当缩放的闭双曲流形）。
\end{proposition}
\begin{proof}
乘积度量的 Ricci 端同态是两因子的直和，故 $\delta_\sigma(g\oplus g_Y)=\max\{\delta_\sigma(g),0\}$。
\end{proof}

log-Weyl 机制本身与维数无关：恒等式~\eqref{eq:logid}对一切 $d\ge3$ 成立，只需一个带无障碍核心修正的 Ricci-flat ALE 空间（预期 $\mathbb C^m/\Z_m$ 上的 Calabi 度量及其余齐性一 Einstein 形变族可以胜任；本报告未验证这一点）。但这些构造对原题的意义受制于如下事实：

\begin{proposition}\status{逻辑事实}
对 $n\ge5$，$Q_n$ 的否定以“存在不承载任何 Einstein 度量的闭 $n$ 维流形”为前提。目前在维数 $\ge5$ 没有任何已知的 Einstein 存在性障碍；“是否每个 $n\ge5$ 维闭流形都承载 Einstein 度量”本身是公开问题。若后者为真，则 $Q_n$ 对一切 $n\ge5$ 平凡成立。
\end{proposition}
乘积只传播输入夹逼，不传播“无 Einstein 度量”：未知的 Einstein 度量不必按乘积分裂。因此高维部分既不能由四维结论推出，也不能用现有工具否定；全维肯定证明则需要在一个符号分支内对一切紧夹逼输入（包括塌缩输入）构造 Einstein 度量，这超出了现有存在性理论（扰动法、Dehn 填充、分歧覆盖、Kähler 方法）的范围。

\section{决定性的开放问题与后续建议}\label{sec:open}

下表列出“解决其中任何一项都会实质性改变原题状态”的问题。它们彼此独立，量词各不相同。

\begin{longtable}{L{0.06\textwidth}L{0.50\textwidth}L{0.36\textwidth}}
\toprule 编号 & 问题 & 对原题的后果\\\midrule\endhead
OP1 & $\CP^2\#\CP^2$ 是否承载 Einstein 度量？（同类：$\#3(S^2\times S^2)$、$2\CP^2\#4\CPb$、$3\CP^2\#3\CPb$；若资料包的 ADE 核延拓通过独立审查，还包括一切 $\#_k(S^2\times S^2)$。）& 任一为“否” $\Rightarrow$ $P_4^+$ 不成立；$P_4^+$ 成立 $\Rightarrow$ 全为“是”。\\
OP2 & 负 Einstein $A_1$ 轨形（如 $(C_g\times C_h)/\Z_2$、含孤立 $\Z_2$ 点的双曲轨形）的混合手性 EH 解消，以及 Dehn-twist 曲面丛 $E_D$，是否承载 Einstein 度量？& 任一为“否” $\Rightarrow$ $P_4^-$ 不成立。\\
OP3 & 问题 W：$T^3\times\R$ 上是否存在两端渐近双曲（共形无穷远为 Yamabe 负环面）的完备 Einstein 度量？& “是” $\Rightarrow$ $P_4^-$ 不成立（定理~\ref{thm:W}）。“否”并配合正则性 $\Rightarrow$ 图流形路线关闭。\\
OP4 & SW 边界带：对 $(\frac13-\gamma)c_1^2(X)\le k\le\frac13c_1^2(X)$ 的固定流形 $X\#k\CPb$，是否有 $e_->0$？相关的正则性假设 (SC)：有界 Ricci、小 $L^2$ 曲率能量的塌缩球上是否存在原度量的统一弱调和展开片？& 某个这样的流形 $e_-=0$ $\Rightarrow$ $P_4^-$ 不成立。(SC) 为“是” $\Rightarrow$ 一致间隙 $\delta_-^2+d\ge c>0$，临界比例附近的部分关闭。\\
OP5 & 对承载正数量曲率度量、满足严格 HT 的四维流形，是否存在新的 Einstein 不存在性障碍？& 否定 $P_4^+$（从而否定 $Q_4$）的必要前提。\\
OP6 & 是否存在不承载 Einstein 度量的闭 $n\ge5$ 维流形？& 否定任何 $Q_n$（$n\ge5$）的必要前提。\\
OP7 & 在一个符号分支内，对一切紧夹逼输入（含塌缩、曲率集中）构造 Einstein 度量。& 证明 $Q_n$ 的唯一途径；四维时已蕴含 OP1 或 OP2 的肯定回答。\\
\bottomrule
\end{longtable}

\paragraph{为什么不能宣称“已解决”。}原题在 $n\ge4$ 的任何完整答案都必须经过上表：肯定答案蕴含 OP1/OP2 的存在性结论（定理~\ref{thm:R}(a)--(c)），否定答案需要 OP5，并在负分支另有突破（OP2、OP3，或 OP4 中边界带上的固定／塌缩序列），高维还受制于 OP6。资料包三百余步的研究与本次核验都没有发现绕过这些问题的途径；本报告给出的归约说明这不是方法选择的问题，而是原题的内在结构。

\paragraph{对后续研究的具体建议。}
\begin{enumerate}
\item \textbf{可独立成文的结果}（均已通过本次核验）：（i）定理~\ref{thm:A1}及推论~\ref{cor:football}、\ref{cor:stability}——对 Ozuch 的 $L^\infty$ 公开问题的肯定回答，含锐首项常数（本报告数值确认）与资料包的通量下界（终审复核了其线性 Wronskian 恒等式与球面情形的 $1/(8L)$）；（ii）定理~\ref{thm:U}、\ref{thm:collapse}——LeBrun 边界的统一间隙与必要塌缩；（iii）定理~\ref{thm:ED}及其二阶、三阶分析；（iv）命题~\ref{prop:morse}、定理~\ref{thm:W}——图流形路线的归约。成文时按第~\ref{sec:verification}节的勘误表修订。
\item \textbf{OP3 的正确数值表述}：以 $T^2$ 对称、带联络扭转的约化系统为第一目标（正交传递类已被完全排除），在共形规范或 Einstein--DeTurck 规范下解二维椭圆边值问题，边界数据为两端的共形无穷远 $[dz^2+e^{v_\infty(z)}dx^2+e^{-v_\infty(z)}dy^2]$ 与常值荷；从带人工源项、可支撑 $\cosh$ 型虫洞的同伦族出发做延拓。Cauchy（打靶）表述因椭圆不适定而不可用。若数值上得到解，再用 Newton--Kantorovich 型后验估计做计算机辅助证明；若延拓分支在源项趋零前总是终止，则寻找相应的 Pohozaev/Noether 型恒等式。
\item \textbf{OP4}：目标应是“原度量”在局部覆盖上的 $C^{1,\alpha}$ 调和片，而不是平滑后度量的控制；可尝试结合 Cheeger--Tian 的塌缩 $\varepsilon$-正则性论证与 Naber--Zhang 的纤维化基本群条件。
\item \textbf{不建议继续投入的方向}：在 SW 临界比例上寻找非塌缩近 Einstein 序列（已被定理~\ref{thm:collapse}排除）；正交传递或固定螺旋标架类的连接器（已被排除）；用 HT 型拓扑筛选寻找正分支反例（命题~\ref{prop:HT}）。
\end{enumerate}

\section{核验记录与勘误建议}\label{sec:verification}

\subsection{范围与方法}
输入：\path{qiushi-engine-yau26.zip}，其 SHA-256 为
\begin{center}\texttt{\small f28ebf9e2f7e0fd450fbefd7346644d054ed58dcef5f82222509f45a1a4d4c44}\end{center}
通读了研究总论、研究地图与六篇研究稿的主体；两份计算附录（约 10\,MB）只按需检索。核验分三层：（1）运行随包程序 \path{scripts/reproduce.py}，五组检查（\texttt{uniform\_gap}、\texttt{weyl10}、\texttt{ed\_ricci}、\texttt{cz\_spike}、\texttt{four\_node}）全部通过；（2）本报告作者从定义出发另写的符号与高精度数值程序（不调用随包代码）；（3）四个相互独立的审稿式核查（统一间隙稿、曲面丛一章、七项定量障碍命题、文献陈述），每个核查都逐式重推并另写程序；（4）对本报告自身新增论断的一次对抗式终审（\path{final_review/}），其指出的措辞过强之处（SW 边界带、monopole 类的表述、SW 窗口推论的适用范围、不可定向底面）均已在正文中更正。

\subsection{独立计算及结果}
\begin{longtable}{L{0.22\textwidth}L{0.48\textwidth}L{0.22\textwidth}}
\toprule 程序 & 内容 & 结果\\\midrule\endhead
\path{v1_core.py} & 一般双轴 Bianchi IX 度量的 Ricci 张量（坐标 Christoffel 符号，符号计算）；标架对角性 & 残差恒为 $0$\\
\path{v1b_family.py} & $q_\Lambda$ 族满足 $\Ric=\Lambda q_\Lambda$（用 ODE 消元） & 三个特征值 $-\Lambda\equiv0$\\
\path{v1c_htilde.py} & $\partial_\Lambda G|_0$、规范场 $V$、$\tilde h$ 的三个分量与渐近式 $H_*$ & 与引理~\ref{lem:core}逐项相等\\
\path{v2_logidentity.py} & 式~\eqref{eq:logid}，$d=4,5$，随机代数 Weyl 张量、任意 $\eta$ & 残差恒为 $0$\\
\path{v3b_football.py} & $S^4/\Z_2$ 及负曲率局部模型上 $g_a$ 的精确 Ricci 特征值，$a=10^{-8},\dots,10^{-128}$，精度 $3\,\mathrm{ex}+60$ 位 & $L\delta\to0.33334$（表~\ref{tab:football}；$a=10^{-128}$ 时 $\delta=0.0046554$）\\
\path{v4_tail.py} & 尾部 $dr^2+e^{2r}(dz^2+e^{v}dx^2+e^{-v}dy^2)$ 的全部 Ricci 缺陷分量 & 与资料包式 \texttt{eq:explicit-tail} 相等\\
\path{audit_bundle/*.py} & 曲面丛 (C1)--(C4)：完整 Ricci 矩阵（两条路线）、双领圈首系数、线性化算子、两个标量 PDE 的数值解、共形恒等式 & 全部吻合；$L\sup\|\Phi\|\to0.985$；$\sup|\Phi(G_\varepsilon)|/\varepsilon^2=0.02053\approx2/\pi^4$\\
\path{audit_obstructions/*.py} & HT、SW 窗口、BCG 阈值、$W$ 与截面曲率不等式、Gaussian Ricci 公式、螺旋标架 ODE、OT 定理的 Bochner 恒等式与热核积分 & 全部吻合；$\epsilon_*=0.2462498$\\
\bottomrule
\end{longtable}

\subsection{审稿结论}
\begin{longtable}{L{0.34\textwidth}L{0.58\textwidth}}
\toprule 对象 & 结论\\\midrule\endhead
统一间隙稿：\texttt{thm:exact}、\texttt{thm:local}、\texttt{thm:proportion}/\texttt{cor:diam}、\texttt{prop:conditional} & 未发现实质缺口。校准恒等式与常数 $109,436,675,42$ 均复核；LeBrun 的负部形式 $\int(\frac23s+2w)_-^2\ge32\pi^2(a^+)^2$ 确为其 Theorem 2.3 证明中的不等式；Cheeger--Tian 无体积下界的 $\varepsilon$-正则性按缩放正确使用；重数传输引理的路径提升论证正确。条件命题的证明为梗概，(SC) 未证（原稿已明示）。\\
曲面丛一章 (C1)--(C4) & 未发现错误。(C1)、(C2) 完全验证；(C3)、(C4) 的全部代数与数值成分验证，所依赖的标准椭圆估计为断言式陈述；一处外部引理（Rupflin--Topping）未能取得原文核对。$e_-(E_D)=0$ 无条件成立。\\
七项定量命题（HT、SW 窗口、BCG、图流形预算、Gaussian 公式与中心收缩、螺旋标架、OT 定理） & 全部正确；需小修处见下。\\
log-Weyl 定理 \texttt{thm:A1} & 正确；本报告第~\ref{sec:logweyl}节给出自足证明并以独立数值确认首项常数。\\
\bottomrule
\end{longtable}

\subsection{勘误与修订建议（供成文时采用）}
\begin{enumerate}
\item 基础篇 \texttt{prop:gap} 的阈值 $1/(1+\sqrt3)$ 有效但非最优；应与长稿 \texttt{FND:HT} 的端点 $1/2$ 统一（见命题~\ref{prop:HT}）。
\item 统一间隙稿：定义“monopole 类”与 $W^+$ 在 $\Lambda^+$ 上的作用约定；注明谱估计是 LeBrun Theorem 2.3 证明中的中间不等式；写明 Cheeger--Tian 定理的缩放（$\Ric=-g$ 到 $|\lambda|\le3$ 规范）；为展开球上的调和半径界补引文（Jost--Karcher/Anderson）；Cheeger--Naber 的式号 (7.12)--(7.14) 未能核对，建议改引其 Definition 2.9 后的注记；“极限已经 Einstein”应改为“极限满足 $\Ric=-h_\infty$”；$s,\delta,p,\eta,A,D$ 等记号重载需清理。
\item 长稿曲面丛一章：双领圈结论的“充分小 $\varepsilon$”实际约为 $\varepsilon\lesssim10^{-3}$（截断带常数约 $14$），宜写明；\texttt{BUN:JQ} 存在前向引用；$u,Q,D,K,a,L$ 记号重载。
\item 全局篇：第 78 行“曲率算子范数”与 Ricci 的 $\|\cdot\|_{\op}$ 冲突，应写 Hilbert--Schmidt（$\|W\|_\infty\ge(1-3\delta)/\sqrt3$ 对谱范数不成立）；Preissmann 步骤需一行引文（$\Z^3$ 边群单射）；\texttt{thm:OT} 的热核部分应引 Bakry--Gentil--Ledoux 或 Erbar--Kuwada--Sturm，并可删去“质量为一”的假设；$\chi_N=2,4,10$ 仅为数值示例。
\item 长稿 \texttt{FND:SWgap}：应显式写出“$\mathcal E\mapsto\mathcal E/(4\pi^2T+\mathcal E/2)$ 在 $T>0$ 时单调”这一步（这是 $K<c$ 的唯一用处）；“$\delta<1/2$”可放宽。
\item Cai--Galloway 定理的衰减条件是对 Ricci 张量（$\Ric=-ng+o(r^2)$）而非数量曲率；$dr^2+\cosh^2r\,g_{T^3}$ 是 $O(r^2)$ 而非 $o(r^2)$ 的例子，与该定理不矛盾。
\item 构造篇引用 Ozuch 的球面商问题时，建议同时引其 2021 年报告幻灯片中明确的表述（“Open question: With $\|\Ric(g_t)-3g_t\|_{L^\infty(g_t)}\to0$?”）；论文正文 \S7.3 的原文本次未能取得。
\end{enumerate}

\subsection{文献核对的边界}
未能核对：Yau 1982 年原书第 26 题的逐字文本；Ozuch（G\&T 2022）\S7.3 的正文；Rupflin--Topping 引理的编号。经检索得到但本报告未读原文的新文献：Li--Liu（\cite{LiLiu}，报道 $(0,\infty)\times T^3$ 上一端 Poincaré--Einstein、另一端 AH 尖点的完备反自对偶 Einstein 度量，共形无穷远为 Yamabe 负）；它与问题 W 相关（说明 Yamabe 负环面无穷远的 Einstein 填充确实存在），但不给出两端外扩的解。检索未发现任何给出 $L^\infty$ 夹逼解消或“对数 Weyl 插值”的已有工作。

\subsection{关于问题 W 的一次探索性数值试验（无结论）}
本次另写了 $T^2$ 对称、Gaussian 规范下的 Einstein 演化程序（\path{wormhole/evolve.py}：$\partial_rq=2qS$，$\partial_rS=q^{-1}\Ric(q)+3I-(\tr S)S$，$z$ 方向 Fourier 拟谱；三维 Ricci 例程通过对角公式与剪切不变性单元测试，误差 $10^{-13}$），并构造了带扭转荷的时间对称喉部初值（Hamilton 约束化为 Ermakov--Pinney 型周期方程，存在性要求 Hill 算子 $-\partial_s^2+\frac94-\frac3{16}v_s^2$ 的最低特征值为负，这与 $\langle e\rangle\ge12$ 一致）。在精确双曲尖点的“波纹切片”上验证表明：约束破坏按最高保留模的速率 $e^{kr/\ell}$ 增长（椭圆 Cauchy 问题不适定），低模截断虽可积分到 $r=6$ 但约束误差为 $O(1)$。因此打靶法不能区分真实奇性与数值失稳，\emph{不支持任何关于问题 W 的结论}；这也解释了资料包中全矩阵试验停滞的原因。建议的替代表述见第~\ref{sec:open}节。

\subsection{产物清单与 SHA-256}
以下文件随报告一并交付（压缩包 \path{yau26_verification_artifacts.zip}；\path{final_review/} 为终审复核脚本）。

{\scriptsize\begin{longtable}{L{0.36\textwidth}L{0.58\textwidth}}
\toprule 文件 & SHA-256\\\midrule\endhead
\path{verify/v1_core.py} & \texttt{510d9e0f094649fd03ee393237ab6acc}\allowbreak\texttt{6899d9f1203970ab9bda3504cd6eb970}\\
\path{verify/v1b_family.py} & \texttt{f54a88095873783294cb36834011266d}\allowbreak\texttt{d5f5a2c0032e2b357cb570caf78a11f7}\\
\path{verify/v1c_htilde.py} & \texttt{9876a694c48dcaf6825612cb43e4c890}\allowbreak\texttt{7d2c3636773587511d750b4ffb0f3705}\\
\path{verify/v2_logidentity.py} & \texttt{665aa565d8f847892c6aaf14334b9c59}\allowbreak\texttt{e7b0b148a0b736b7b02314537dc7871d}\\
\path{verify/v3_football.py} & \texttt{e9d7fb50196ea88fd970eeb8785cc8b0}\allowbreak\texttt{23895a79337685d165510189605313f1}\\
\path{verify/v3b_football.py} & \texttt{cf43e43b221c0181a5458d73abd17abd}\allowbreak\texttt{9f8b0f03bf3a05c5dac09aef495ed33a}\\
\path{verify/v4_tail.py} & \texttt{71b51885c5e400bb1465f6ae7a418a63}\allowbreak\texttt{dbdf4525b1eb1fbfbed32a400157685f}\\
\path{verify/football_1e128.log} & \texttt{9765f19bc3334c11979f9520f26dc079}\allowbreak\texttt{a36284a3fafdd4f59f941568444209d9}\\
\path{verify/football_minus.log} & \texttt{1aabb205412190285f2112c77295af6b}\allowbreak\texttt{81668e733782b8703efab837eaa97008}\\
\path{verify/football_plus.log} & \texttt{071a7c639a6c08e4f7d9b6193104ff28}\allowbreak\texttt{d6e7a2da89105ec72991f7a347826e10}\\
\path{wormhole/evolve.py} & \texttt{de2595f3bf44aa85207ca2fb0e662dbe}\allowbreak\texttt{170aa4d9fa28cc6ef99d6526866497a0}\\
\path{wormhole/ot_test.py} & \texttt{ebb1ac7f7f7c61f786c2ab4608f38350}\allowbreak\texttt{ae5b847e00948950edc0517eccfec96c}\\
\path{wormhole/throat.py} & \texttt{4a11c611f6bc2935b169cdce156e7969}\allowbreak\texttt{406da127311e235b845a4fa18f75021a}\\
\path{wormhole/validate.py} & \texttt{76f8bc9c7043263cd2a3e84b06188ae6}\allowbreak\texttt{90b574c4311509bdbae57888fd05b1ab}\\
\path{audit_bundle/c1_explicit_crosscheck.py} & \texttt{e9d4794256c154add18f4c86e45dc979}\allowbreak\texttt{a9cc6214a7e94c011c90132a4d75c0ff}\\
\path{audit_bundle/c1_ricci_slow.py} & \texttt{c8a68904344408ef72de332390ce209d}\allowbreak\texttt{a79dc034c99e697e50048bb4de41c051}\\
\path{audit_bundle/c2_double_collar.py} & \texttt{6dda30e8aeeabe21d15e104b02c62fbc}\allowbreak\texttt{f80b9a99cea1109f4c8ecf7658b227d5}\\
\path{audit_bundle/c3_linearization.py} & \texttt{65fa407c07b72188be9e7e5d364f6ed5}\allowbreak\texttt{f5e539030b3ba7163bad356ab1fb1ee1}\\
\path{audit_bundle/c3_nonlinear_residual.py} & \texttt{658062b9124ad6b1eb8a33d412634a96}\allowbreak\texttt{aae3b1d154e55afd962d6912bd405a89}\\
\path{audit_bundle/c3_pde_numeric.py} & \texttt{a0e04e05d81ff10e06ed7c621bb24758}\allowbreak\texttt{7b8997f304516ac79181dbfeb317734f}\\
\path{audit_bundle/c4_JQ_theta_dependent.py} & \texttt{8698392e44c968804e565b467600a762}\allowbreak\texttt{64a702829322e59e04df0859915ce0c0}\\
\path{audit_bundle/c4_conformal_identity.py} & \texttt{a871d838a3eb751a68485e9d91b2e733}\allowbreak\texttt{0ee48c14082d2c1ab2ceae987ffa9de0}\\
\path{audit_bundle/c4_pairings.py} & \texttt{a1c588259daabab47bae5250bbfff7ab}\allowbreak\texttt{e5617e5590b79b529df23cc856925d65}\\
\path{audit_obstructions/itemABC_algebra.py} & \texttt{7fb92f6514c922788f3e700fe50a210c}\allowbreak\texttt{9abdd25ca8ffce338fe01394e0001c05}\\
\path{audit_obstructions/itemD_weyl_sectional.py} & \texttt{4464b4e92259abf7a460fc72ee187161}\allowbreak\texttt{67a8677e36b3b375c11519eec5279e3b}\\
\path{audit_obstructions/itemE_gaussian.py} & \texttt{3000e4f40ee6ad9724f970ce441e37a3}\allowbreak\texttt{c8d205382ebc83bb3d8d92c667d555bc}\\
\path{audit_obstructions/itemE_jets_numeric.py} & \texttt{c847ed254e5d95e7dabd4ea2c3ba829f}\allowbreak\texttt{e639372c12d7881c138cbf56df9236c6}\\
\path{audit_obstructions/itemE_jets_numeric_lib.py} & \texttt{69cf97c4b408bdddc7348051a63fd252}\allowbreak\texttt{a897da62b7051e112c48b2aaaa22f5ea}\\
\path{audit_obstructions/itemF_helical.py} & \texttt{af4f5f36a8f596c4e8cb36d66e8b0c1f}\allowbreak\texttt{f25bb3a62d4236feb25cb93a68bacaec}\\
\path{audit_obstructions/itemG_OT.py} & \texttt{c9549d0bb729d14700f62d4273f5807a}\allowbreak\texttt{1588c016fc073e1075ed44bc3dd467cc}\\
\path{audit_obstructions/itemG_bochner_identity.py} & \texttt{d1eba5606d2d8e632c1820e44db05eb3}\allowbreak\texttt{37bf08c0c8ad1ff1c12862877b9b6413}\\
\path{audit_obstructions/ricci_tools.py} & \texttt{fbf60942df2bb0e785ee9b1870520158}\allowbreak\texttt{a2b068d95e698cd0757c474cb50e6a82}\\
\path{final_review/r1_weylgauge.py} & \texttt{35b5bab5632a90e993928997e3239b46}\allowbreak\texttt{a5d4e1ddff5e6244d5afc8be78a0719a}\\
\path{final_review/r2_annulus_generic.py} & \texttt{583d9831ec6797859cd936a13e4af2b4}\allowbreak\texttt{f410de6607d97b3a15a81b11608c9d96}\\
\path{final_review/r2_core.py} & \texttt{2295ee873e23f5d13717fff9ba5fb8db}\allowbreak\texttt{6855ac88071c5706525b79e29b4f0beb}\\
\path{final_review/r2_core_region.py} & \texttt{8170e4fd9941f1784ba4fbcd740e0fd1}\allowbreak\texttt{fc663e2db468c1325b8cb42b13a5f4d3}\\
\path{final_review/r2_table_check.py} & \texttt{1edb625b5b3ae2bf391532b6fea121cb}\allowbreak\texttt{53c84df2dcacfa3281e8fc7f868ec239}\\
\path{final_review/r4_OT_identities.py} & \texttt{686dba8ba0a18efd107dfbdbba50ed92}\allowbreak\texttt{d744a8ef1f9828ab7f1bf5e6b05d9586}\\
\path{final_review/r6_gaussian_blocks.py} & \texttt{49de53469175c021ed95da3d909f943a}\allowbreak\texttt{cbb41129e9d1a1254e3c3b6e9a1e6125}\\
\path{final_review/r6_tail_interp.py} & \texttt{0bc54367a66051c2042f5630aa8f6836}\allowbreak\texttt{d5d0310820f8a549225156f16116603c}\\
\path{final_review/r7_SWwindow.py} & \texttt{bf540c495f3c02a02763b691c2867eb9}\allowbreak\texttt{603fa175d62e6a5ae6e54e06b0ced736}\\
\path{final_review/r7_bundle_check.py} & \texttt{53d235bd77cb924c0f0a0c595af282a0}\allowbreak\texttt{481f3fbe2ca1788c377ccb4b95fb9584}\\
\path{repro_run1/cz_spike.json} & \texttt{1df305a78321769657f9a425761c79e8}\allowbreak\texttt{a783e8971d5c39e2366682ae42144137}\\
\path{repro_run1/ed_ricci.json} & \texttt{03e7ea671a52b1d265f88005769e5325}\allowbreak\texttt{a7e24f7ebc9fec87f6a613b4b062df23}\\
\path{repro_run1/four_node.json} & \texttt{6e9e4e96095ddd2d791ef0a3075c7666}\allowbreak\texttt{33179a57fe0d26e29262659aa0a44943}\\
\path{repro_run1/manifest.json} & \texttt{a32274099b3704c0dd5687e19ae424e3}\allowbreak\texttt{292ebc186d8660db26a02cdb58db6f26}\\
\path{repro_run1/uniform_gap.json} & \texttt{3d38dcdea1299af1f8abf1aed8bef420}\allowbreak\texttt{8c2d114609dd63b5ba5a97f8d78ed100}\\
\path{repro_run1/weyl10.json} & \texttt{fc42e4dd942ea007a5888318445dc78b}\allowbreak\texttt{ab873be369afbe35263b66026577d36f}\\
\bottomrule\end{longtable}}


\begin{thebibliography}{99}\small
\bibitem{Yau} S.-T. Yau, \emph{Problem section}, Seminar on Differential Geometry, Ann. of Math. Stud. 102, Princeton Univ. Press, 1982, 669--706.
\bibitem{Pkg} 求是引擎资料包 \texttt{qiushi-yau26}：\path{research/20260910.tex}（基础篇）、\path{20260912.tex}（长稿）、\path{20260916_uniform_gap.tex}、\path{20260916_rigidity.tex}、\path{20260916_construction.tex}、\path{20260916_global.tex}，及 \path{references/}、\path{scripts/}。
\bibitem{Hamilton} R. Hamilton, Three-manifolds with positive Ricci curvature, J. Differential Geom. 17 (1982), 255--306.
\bibitem{Perelman} G. Perelman, Ricci flow with surgery on three-manifolds, arXiv:math/0303109；B. Kleiner, J. Lott, Notes on Perelman's papers, Geom. Topol. 12 (2008).
\bibitem{Lohkamp} J. Lohkamp, Metrics of negative Ricci curvature, Ann. of Math. 140 (1994), 655--683.
\bibitem{Anderson89} M. Anderson, Ricci curvature bounds and Einstein metrics on compact manifolds, J. Amer. Math. Soc. 2 (1989), 455--490.
\bibitem{Biquard} O. Biquard, Désingularisation de métriques d'Einstein I, Invent. Math. 192 (2013), 197--252.
\bibitem{OzuchII} T. Ozuch, Noncollapsed degeneration of Einstein 4-manifolds I, II, Geom. Topol. 26 (2022), 1483--1528, 1529--1634（arXiv:1909.12960）；报告幻灯片 “Noncollapsed degeneration of Einstein 4-manifolds”，2021。
\bibitem{OzuchCPAM} T. Ozuch, Integrability of Einstein deformations and desingularizations, Comm. Pure Appl. Math. 77 (2024)（arXiv:2105.13193）。
\bibitem{HuangOzuch} Y. Huang, T. Ozuch, Regularity of Einstein 5-manifolds via 4-dimensional gap theorems, arXiv:2512.21317。
\bibitem{BK} S. Brendle, N. Kapouleas, Gluing Eguchi--Hanson metrics and a question of Page, Comm. Pure Appl. Math. 70 (2017), 1366--1401.
\bibitem{OR} P. Orlik, F. Raymond, Actions of the torus on 4-manifolds I, Trans. Amer. Math. Soc. 152 (1970), 531--559.
\bibitem{Bonahon} F. Bonahon, Difféotopies des espaces lenticulaires, Topology 22 (1983), 305--314.
\bibitem{LeBrun01} C. LeBrun, Ricci curvature, minimal volumes, and Seiberg--Witten theory, Invent. Math. 145 (2001), 279--316.
\bibitem{LeBrun03} C. LeBrun, Einstein metrics, four-manifolds, and differential topology, Surveys in Differential Geometry VIII (2003), 235--255.
\bibitem{LeBrun09} C. LeBrun, Einstein metrics, complex surfaces, and symplectic 4-manifolds, Math. Proc. Cambridge Philos. Soc. 147 (2009), 1--8.
\bibitem{IL} M. Ishida, C. LeBrun, Curvature, connected sums, and Seiberg--Witten theory, Comm. Anal. Geom. 11 (2003), 809--836.
\bibitem{Viaclovsky} J. Viaclovsky, Critical metrics for Riemannian curvature functionals, IAS/Park City Math. Ser. 22 (2016).
\bibitem{Derd} A. Derdziński, Self-dual Kähler manifolds and Einstein manifolds of dimension four, Compositio Math. 49 (1983), 405--433.
\bibitem{CT} J. Cheeger, G. Tian, Curvature and injectivity radius estimates for Einstein 4-manifolds, J. Amer. Math. Soc. 19 (2006), 487--525.
\bibitem{CN} J. Cheeger, A. Naber, Regularity of Einstein manifolds and the codimension 4 conjecture, Ann. of Math. 182 (2015), 1093--1165.
\bibitem{Li} Y. Li, Smoothing Riemannian metrics with bounded Ricci curvatures in dimension four, I, II, Adv. Math. 223 (2010)；Ann. Global Anal. Geom. 41 (2012)。
\bibitem{BCG} G. Besson, G. Courtois, S. Gallot, Entropies et rigidités des espaces localement symétriques de courbure strictement négative, Geom. Funct. Anal. 5 (1995), 731--799.
\bibitem{Sambusetti} A. Sambusetti, An obstruction to the existence of Einstein metrics on 4-manifolds, Math. Ann. 311 (1998), 533--547.
\bibitem{DiCerbo} L. F. Di Cerbo, Extended graph 4-manifolds, and Einstein metrics, Ann. Math. Qué. 48 (2024), 269--276（arXiv:2106.13279）。
\bibitem{CSS} C. Connell, P. Suárez-Serrato, On higher graph manifolds, Int. Math. Res. Not. 2019, 1281--1311.
\bibitem{CG} M. Cai, G. Galloway, Boundaries of zero scalar curvature in the AdS/CFT correspondence, Adv. Theor. Math. Phys. 3 (1999), 1769--1783.
\bibitem{LiLiu} Li, Liu, arXiv:2605.25462（经检索得到，本报告未读原文）。
\end{thebibliography}

\end{document}
